\documentclass[11pt,a4paper]{article}
\usepackage[utf8]{inputenc}
\DeclareUnicodeCharacter{03C4}{\ensuremath{\tau}}
\usepackage[T1]{fontenc}
\usepackage[english]{babel}
\usepackage{amsmath,amssymb,amsthm}
\usepackage{xcolor}
\usepackage{forest}
\usepackage{booktabs}
\usepackage{enumitem}
\usepackage{longtable}
\usepackage[margin=2.6cm]{geometry}
\usepackage[colorlinks=true,linkcolor=blue!60!black,citecolor=green!40!black]{hyperref}

\title{Multirate infinitesimal step methods with polynomial coupling:\\
frozen formulations, relations to other method classes, and order conditions}
\author{Oswald Knoth}
\date{Draft, \today}

% ---------------------------------------------------------------------
\newtheorem{theorem}{Theorem}[section]
\newtheorem{proposition}[theorem]{Proposition}
\newtheorem{lemma}[theorem]{Lemma}
\newtheorem{conjecture}[theorem]{Conjecture}
\theoremstyle{definition}
\newtheorem{definition}[theorem]{Definition}
\newtheorem{example}[theorem]{Example}
\newtheorem{question}{Question}
\theoremstyle{remark}
\newtheorem{remark}[theorem]{Remark}

\newcommand{\R}{\mathbb{R}}
\newcommand{\T}{\mathcal{T}}
\newcommand{\dd}{\,\mathrm{d}}
\newcommand{\todo}[1]{\textcolor{red}{[\textbf{TODO:} #1]}}

% Bicoloured trees -----------------------------------------------------
\forestset{
  b/.style={circle,draw=black,fill=black,inner sep=1.3pt},
  w/.style={circle,draw=black,fill=white,inner sep=1.3pt},
  bitree/.style={for tree={grow=90,s sep=3pt,l=5pt,l sep=0pt,
                 parent anchor=center,child anchor=center,edge={thin}},
                 baseline}
}
\newcommand{\btree}[1]{\begin{forest}bitree #1\end{forest}}
% order 1
\newcommand{\Tb}{\btree{[,b]}}
% order 2
\newcommand{\Tbb}{\btree{[,b [,b]]}}
\newcommand{\Tbw}{\btree{[,b [,w]]}}
% order 3
\newcommand{\Tbbb}{\btree{[,b [,b][,b]]}}
\newcommand{\Tbbw}{\btree{[,b [,b][,w]]}}
\newcommand{\Tbww}{\btree{[,b [,w][,w]]}}
\newcommand{\Tbbbl}{\btree{[,b [,b [,b]]]}}
\newcommand{\Tbbwl}{\btree{[,b [,b [,w]]]}}
\newcommand{\Tbwbl}{\btree{[,b [,w [,b]]]}}
\newcommand{\Tbwwl}{\btree{[,b [,w [,w]]]}}
% selected order 4
\newcommand{\Tfourbbw}{\btree{[,b [,b][,b][,w]]}}
\newcommand{\Tfourbww}{\btree{[,b [,b][,w][,w]]}}
\newcommand{\Tfourbbl}{\btree{[,b [,b [,b]][,w]]}}
\newcommand{\Tfourbbwl}{\btree{[,b [,b [,b][,w]]]}}

\begin{document}
\maketitle

\begin{abstract}
Multirate infinitesimal step (MIS) methods integrate a slow part of an
ordinary differential equation with a Runge--Kutta-like outer scheme and
the fast part by the solution of auxiliary (``infinitesimal'') initial
value problems. They are the mathematical form of the split-explicit time
integrators used in numerical weather prediction.

We formulate MIS methods for $y'=F(y,y)$, where the first argument is
frozen at stage values and the second argument is integrated in time.
Both the coupling to the frozen right-hand sides and the coupling to
previous stage increments are polynomials in the inner time. For the
compressible Euler equations we discuss the partitioning in potential
temperature, internal and total energy form, and two ways of freezing
the pressure in the potential temperature formulation.

Several known method classes appear as special cases: classical MIS and
split-explicit Runge--Kutta schemes, MRI-GARK methods, a class of
exponential Runge--Kutta methods, and commutator-free Lie group methods.
With bicoloured B-series whose coefficients are polynomials in the inner
time, we derive the order conditions in a Runge--Kutta-like form and give
them explicitly up to order 4. For methods that are exact for the fast
part we prove that order 3 requires only the base method conditions and
one coupling condition, and we count the coupling conditions up to order
5. We further discuss stiff order conditions, spatial discretizations
(finite volumes, DG with the LMARS flux), the solution and cost of the
linear inner problems, and the required accuracy of inexact inner
solutions.
\end{abstract}

\tableofcontents

% =====================================================================
% =====================================================================
\section{Introduction}
\label{sec:motivation}

\subsection{Motivation}

Atmospheric models for numerical weather prediction (NWP) and climate,
and many CFD codes for low-Mach-number compressible flow, solve the
compressible Euler (or Navier--Stokes) equations. The solution contains
processes on very different time scales:
\begin{itemize}
\item slow: advection by the flow (wind speeds up to $\approx100\,$m/s)
  and the physical parameterizations;
\item fast: sound waves ($\approx340\,$m/s) and, in high resolution,
  gravity waves.
\end{itemize}
Sound waves carry almost no meteorologically relevant energy, but an
explicit scheme has to resolve them in time. Two remedies are in use:
\begin{itemize}
\item (semi-)implicit treatment of the fast terms, horizontally
  explicit/vertically implicit (HEVI) splittings, IMEX schemes
  \cite{BoscarinoFilbetRusso2016,GiraldoKellyConstantinescu2013};
\item \emph{split-explicit} schemes: the slow tendencies are evaluated
  with a large step $h$ and held fixed (or interpolated) while the fast
  terms are integrated with many small steps
  \cite{KlempWilhelmson1978,WickerSkamarock2002,KlempSkamarockDudhia2007}.
  The Runge--Kutta-based split-explicit scheme of Wicker and Skamarock is
  used for example in WRF and in COSMO \cite{Baldauf2010}.
\end{itemize}
The \emph{multirate infinitesimal step} (MIS) methods
\cite{KnothWolke1998,WenschKnothGalant2009,KnothWensch2014} provide the
mathematical framework for split-explicit schemes. In the analysis, the
fast part is assumed to be integrated exactly (``infinitesimally''). The
analysis then yields order conditions for the outer method. In practice
the fast part is integrated by a few small steps of a cheap explicit
method, or partly implicitly.

\subsection{Partitioning by arguments}

Split-explicit schemes do not split the right-hand side into a sum of a
slow and a fast term. Instead, \emph{slowly varying quantities are
frozen} at the values of a previous stage, while the remaining
dependence on the solution is integrated in time with small steps. This
is formalized by writing the (semi-discrete) equation $y'=\mathcal F_h(y)$
as
\begin{equation}\label{eq:G}
y' = F(y,y), \qquad F:\R^N\times\R^N\to\R^N,\qquad F(y,y)=\mathcal F_h(y), \tag{$\mathcal G$}
\end{equation}
where the \emph{first argument} contains the slow (frozen) and the
\emph{second argument} the fast (integrated) dependence. In a method the
first argument is set to a stage value $Y_j$ and the second argument is
the solution $Z$ of an inner initial value problem
$Z'=F(Y_j,Z)$. The partitioning is a modelling decision. It is not
unique, and its quality decides how large the outer step can be.
Throughout, the index $y$ ($z$) at a quantity means evaluation at the
first (second) argument.

\paragraph{Euler equations.}
In two space dimensions $(x,z)$ with potential temperature $\theta$ the
compressible Euler equations read
\begin{equation}\label{eq:euler}
\begin{aligned}
 \partial_t\rho&=-\partial_x(\rho u)-\partial_z(\rho w),\\
 \partial_t(\rho u)&=-\partial_x(\rho uu)-\partial_z(\rho wu)-\partial_x p,\\
 \partial_t(\rho w)&=-\partial_x(\rho uw)-\partial_z(\rho ww)-\partial_z p-\rho g,\\
 \partial_t(\rho\theta)&=-\partial_x(\rho u\theta)-\partial_z(\rho w\theta),
\end{aligned}
\end{equation}
with the equation of state and the Exner pressure $\Pi$
\begin{equation*}
 p=p_0\Big(\frac{R_d\,\rho\theta}{p_0}\Big)^{\gamma_a},\qquad
 \Pi=\Big(\frac p{p_0}\Big)^{R_d/c_p},\qquad
 p=R_d\,\Pi\,\rho\theta,\qquad
 \frac{\partial p}{\partial(\rho\theta)}=\gamma_a\frac{p}{\rho\theta}
 =\gamma_aR_d\Pi=\frac{c^2}{\theta},
\end{equation*}
$\gamma_a=c_p/c_v$, $c^2=\gamma_ap/\rho$. The state is
$y=(\rho,\rho u,\rho w,\rho\theta)$; $\mathbf v=(u,w)$ and $\theta$ are
derived quantities. The fast (acoustic) subsystem consists of
\begin{itemize}
\item the mass flux divergence in the equations for $\rho$ and
  $\rho\theta$,
\item the pressure gradient,
\item and the buoyancy term.
\end{itemize}
The advecting velocities in the momentum equations and the advected
potential temperature change slowly. Marking these frozen quantities in
\textcolor{red}{red}, the partitioned system is
\begin{equation}\label{eq:eulerred}
\begin{aligned}
 \partial_t\rho&=-\partial_x(\rho u)-\partial_z(\rho w),\\
 \partial_t(\rho u)&=\textcolor{red}{-\partial_x(\rho uu)-\partial_z(\rho wu)}-\partial_x p,\\
 \partial_t(\rho w)&=\textcolor{red}{-\partial_x(\rho uw)-\partial_z(\rho ww)}-\partial_z p-\rho g,\\
 \partial_t(\rho\theta)&=-\partial_x(\rho u\textcolor{red}{\theta})-\partial_z(\rho w\textcolor{red}{\theta}).
\end{aligned}
\end{equation}
Red quantities are evaluated at the first argument $y$, all others at
the second argument $z$. The mass flux in the $\rho\theta$ equation is
fast, the transported $\theta$ is frozen:
$-\nabla\cdot\big((\rho\mathbf v)_z\theta_y\big)$. The pressure
gradient still has to be partitioned. Since $p$ is a nonlinear function
of $\rho\theta$, there are several possibilities (Section~\ref{sec:frozen}).

\paragraph{Problem classes.}
Important special cases of \eqref{eq:G} are
\begin{align}
 &\text{affine in the fast argument:} && F(y,z)=f(y)+\mathfrak B\big(w(y)\big)z, \tag{$\mathcal B$}\label{eq:B}\\
 &\text{additive splitting:} && F(y,z)=f(y)+g(z), \tag{$\mathcal A$}\label{eq:A}\\
 &\text{constant linear fast part:}   && F(y,z)=f(y)+Lz, \tag{$\mathcal L$}\label{eq:L}\\
 &\text{frame representation:} && F(y,z)=\textstyle\sum_{k=1}^K f_k(y)\,E_k(z). \tag{$\mathcal F$}\label{eq:F}
\end{align}
In \eqref{eq:B} the linear operator $\mathfrak B(w)$ depends linearly on
a set of frozen coefficient fields $w=w(y)$. This is the structure of
the frozen Euler equations (Section~\ref{sec:bilinear}).
Every \eqref{eq:G} system can be written as a partitioned system by
doubling the unknowns, $u'=F(u,v)$, $v'=F(u,v)$, with $u(0)=v(0)$. Hence
partitioned (P-)series apply. Since both components are equal, the
colour of the root plays no role (Section~\ref{sec:trees}).

\subsection{Other energy variables: a proposal for the partitioning}
\label{sec:energy}

Many CFD codes, and some atmospheric models, use internal or total
energy instead of $\rho\theta$. The same principle applies: transported
specific quantities and advecting velocities are frozen, mass fluxes and
the pressure are fast. We propose the following partitionings; in the
rest of the paper we use the potential temperature formulation
\eqref{eq:eulerred}, as WRF and COSMO do.

\paragraph{Internal energy.}
With $\rho e$, $p=(\gamma_a-1)\rho e$ and the enthalpy $h=e+p/\rho$, the
energy equation $\partial_t(\rho e)=-\nabla\cdot(\rho e\mathbf v)-p\nabla\cdot\mathbf v$
is written as $-\nabla\cdot(\rho\mathbf v\,h)+\mathbf v\cdot\nabla p$
and partitioned as
\begin{equation*}
 \partial_t(\rho e)=-\nabla\cdot\big((\rho\mathbf v)_z\,\textcolor{red}{h}\big)
   +\textcolor{red}{\mathbf v}\cdot\nabla p_z,
 \qquad p_z=(\gamma_a-1)(\rho e)_z .
\end{equation*}
\begin{itemize}
\item The pressure is linear in the state, so no freezing of the
  equation of state is needed in the momentum equations.
\item The fast system has $c_f^2=(\gamma_a-1)h=\gamma_ap/\rho=c^2$, the
  full sound speed.
\item $F(y,y)$ is the original right-hand side. The term
  $\mathbf v\cdot\nabla p$ is non-conservative.
\end{itemize}

\paragraph{Total energy.}
With $\rho E=\rho e+\tfrac12\rho|\mathbf v|^2$, the total enthalpy
$H=E+p/\rho$ and $p=(\gamma_a-1)(\rho E-\tfrac12\rho|\mathbf v|^2)$:
\begin{gather*}
 \partial_t(\rho E)=-\nabla\cdot\big((\rho\mathbf v)_z\,\textcolor{red}{H}\big)-g(\rho w)_z,\\
 p_z=(\gamma_a-1)\big((\rho E)_z-\tfrac12\textcolor{red}{\mathbf v}\cdot(\rho\mathbf v)_z\big).
\end{gather*}
\begin{itemize}
\item The pressure is bilinear in the frozen velocity and the fast
  momentum.
\item $F(y,y)$ is the original right-hand side, and the formulation is
  conservative.
\item The fast sound speed is $c_f^2=(\gamma_a-1)H$. It equals $c^2$ for
  $\mathbf v=0$ and exceeds it by $(\gamma_a-1)|\mathbf v|^2/2$ otherwise.
\end{itemize}

In both energy formulations the pressure is linear or bilinear in the
conserved variables, so its partitioning is unambiguous. Both lead to the
structure \eqref{eq:B}. In the potential temperature formulation, $p$ is
a nonlinear function of $\rho\theta$. This leads to the options of
Section~\ref{sec:frozen}.

\subsection{Outline}

\begin{itemize}
\item Section~\ref{sec:frozen}: partitioning of the pressure and the
  bilinear structure.
\item Sections~\ref{sec:methods}--\ref{sec:exponential}: the MIS methods
  with polynomial coupling, their relation to other method classes, the
  B-series, and the order conditions, explicitly up to order 4.
\item Section~\ref{sec:discretization}: spatial discretizations (finite
  volumes, DG with LMARS).
\item Section~\ref{sec:inner}: the solution of the inner problems and
  their cost.
\item Section~\ref{sec:nonexact}: inexact inner solutions.
\end{itemize}

\section{Frozen formulations of the compressible Euler equations}
\label{sec:frozen}

\subsection{Frozen terms common to all options}

As in \eqref{eq:eulerred}:
\begin{itemize}
\item the mass equation and the buoyancy $-\rho g$ are taken in the fast
  argument;
\item in the $\rho\theta$ equation the mass flux is fast and $\theta$ is
  frozen;
\item momentum advection is frozen completely or, as a variant, written
  as $-\nabla\cdot\big((\rho\mathbf v)_z\otimes\mathbf v_y\big)$.
\end{itemize}
The options differ in the treatment of the pressure.

\subsection{Option A: differentiate the pressure, freeze $\partial p/\partial(\rho\theta)$}

The pressure gradient is written by the chain rule as
$\nabla p=\frac{\partial p}{\partial(\rho\theta)}\nabla(\rho\theta)$ and
the coefficient is frozen:
\begin{equation*}
 -\nabla p\ \leadsto\ -a(y)\,\nabla(\rho\theta)_z,\qquad
 a=\frac{\partial p}{\partial(\rho\theta)}=\frac{c^2}{\theta}.
\end{equation*}
This is the linearized form of the original notes. The fast system
carries the full sound speed $c$, and the frozen coefficient carries no
acoustic information.

On the continuous level $F_A(y,y)$ is the original right-hand side. On
the discrete level, however,
$a(y)\nabla_h(\rho\theta)\ne\nabla_hp(y)$ (the difference is of the size
of the truncation error). In particular, a discretely hydrostatically
balanced state, $\partial_{z,h}p=-\rho g$, is \emph{not} a steady state of
$F_A(y,y)$. This is a known source of spurious circulations over steep
terrain. Remedies are a reference-state formulation
($p=\bar p(z)+p'$, linearizing only $p'$) or the following consistent
variant.

\paragraph{Option A$'$: consistent linearization about the frozen state.}
Let $\mathfrak A(y)$ be the linear operator obtained by freezing as in
Option~A (fast mass flux, frozen $\theta$, frozen $a$). Define
\begin{equation}\label{eq:optionAprime}
 F_{A'}(y,z)=\mathcal F_h(y)+\mathfrak A(y)\,(z-y).
\end{equation}
Then $F_{A'}(y,y)=\mathcal F_h(y)$ exactly: the original discretization
and its discrete balances are preserved. Moreover $\partial_zF_{A'}=\mathfrak A(y)$,
and the remaining slow part is
\begin{equation*}
 \partial_yF_{A'}(y,y)=\mathcal F_h'(y)-\mathfrak A(y).
\end{equation*}
In the pressure term this difference is
$-\nabla_h(a\,\delta)+a\nabla_h\delta=-\delta\,\nabla_ha+O(\Delta x)$.
It involves only the gradient of the background coefficient, a term of
size $c^2/(\theta H)$ with scale height $H$, comparable to buoyancy
frequencies and thus slow. The construction is the
MIS analogue of linearizing the acoustic step about the current large
time step state as in \cite{KlempSkamarockDudhia2007}, and of
Rosenbrock-type linearizations. Since $\mathfrak A(y)$ is evaluated at
every stage, it is not a $W$-method.

\subsection{Option B: split $p=\frac{p}{\rho\theta}\cdot\rho\theta$, freeze the first factor}

Write $p=\pi\cdot\rho\theta$ with $\pi=p/(\rho\theta)=R_d\Pi$ and freeze
$\pi$:
\begin{equation*}
 -\nabla p\ \leadsto\ -\nabla\big(\pi(y)\,(\rho\theta)_z\big).
\end{equation*}
\begin{itemize}
\item $F_B(y,y)=\mathcal F_h(y)$ holds exactly on the discrete level, for
  any discretization of the gradient, since $\pi_y(\rho\theta)_y=p(y)$
  cellwise. Discrete hydrostatic balance is preserved, and no extra
  defect term as in \eqref{eq:optionAprime} is needed.
\item The fast part sees only
  $\partial_z(\pi\,\rho\theta)=\pi=c^2/(\gamma_a\theta)$. Hence the fast
  system has the reduced sound speed $c/\sqrt{\gamma_a}\approx0.85\,c$.
\item Since $\partial\pi/\partial(\rho\theta)=(\gamma_a-1)\pi/(\rho\theta)$,
  the slow part contains the remaining fraction
  $(\gamma_a-1)/\gamma_a\approx0.29$ of the acoustic coupling
  $\nabla(a\,\delta\rho\theta)$. The outer (explicit) method therefore
  sees an acoustic component. It is not clear a priori whether this
  restricts the large step $h$, because the fast and slow
  parts of one wave equation interact. This has to be checked with the
  linear stability function of a split acoustic test problem
  (Section~\ref{sec:open}).
\end{itemize}

\subsection{Comparison}

\begin{center}
\begin{tabular}{lccc}
\toprule
 & A & A$'$ & B\\
\midrule
$F(y,y)=\mathcal F_h(y)$ (discrete) & no & yes & yes\\
discrete hydrostatic balance preserved & no$^{*}$ & yes & yes\\
sound speed in fast part & $c$ & $c$ & $c/\sqrt{\gamma_a}$\\
acoustic coupling in slow part & none & none ($-\delta\nabla_h a$) & $(\gamma_a-1)/\gamma_a$\\
frozen fields per stage & $\theta$, $a$ & $\theta$, $a$, $\mathfrak A(Y_j)Y_j$ & $\theta$, $\pi$\\
inner-step operator & \multicolumn{3}{c}{same bilinear structure \eqref{eq:B}}\\
\bottomrule
\end{tabular}

\smallskip
{\small $^*$ unless a hydrostatic reference state is subtracted.}
\end{center}

\subsection{Bilinear structure and summation of the frozen terms}
\label{sec:bilinear}

In all three options the frozen right-hand side is affine in $z$, and
the operator acting on $z$ depends \emph{linearly} on frozen coefficient
fields $w(y)$:
\begin{equation}\label{eq:bilinear}
 F(y,z)=f(y)+\mathfrak B_0z+\mathfrak B_1\big(w(y)\big)z ,
\end{equation}
with
\begin{itemize}
\item $\mathfrak B_0$: the $w$-independent linear terms (mass flux
  divergence, buoyancy);
\item $(w,z)\mapsto\mathfrak B_1(w)z$: bilinear. Examples are
  $\theta_y(\rho\mathbf v)_z$, $a_y\nabla(\rho\theta)_z$,
  $\nabla(\pi_y(\rho\theta)_z)$, and $(\rho\mathbf v)_z\otimes\mathbf v_y$
  in the variant.
\end{itemize}
For each option:
\begin{itemize}
\item Option A: $w=(\theta,a)$ and $f(y)$ is the frozen momentum
  advection.
\item Option B: $w=(\theta,\pi)$ and $f(y)$ is the frozen momentum
  advection.
\item Option A$'$: $f(y)=\mathcal F_h(y)-\mathfrak B_0y-\mathfrak B_1(w(y))y$.
\end{itemize}

\begin{lemma}[summation of frozen terms]\label{lem:summation}
For \eqref{eq:bilinear} and any polynomials $p_{ij}$,
\begin{gather*}
 \sum_{j<i}p_{ij}(\tau)F(Y_j,z)
 =\sum_{j<i}p_{ij}(\tau)f(Y_j)+q_i(\tau)\mathfrak B_0z
  +\mathfrak B_1\big(W_i(\tau)\big)z,\\
 W_i(\tau)=\sum_{k=0}^{K}\tau^kW_{ik},\qquad W_{ik}=\sum_{j<i}a_{ijk}\,w(Y_j),
\end{gather*}
with $q_i=\sum_jp_{ij}$.
\end{lemma}

Thus each inner step needs \emph{one} application of the bilinear
operator, independent of the number of stages $j<i$ that are coupled.
This holds also for polynomial coupling: the $K+1$ combined fields
$W_{ik}$ are formed once per stage, and $W_i(\tau)$ is formed at each
inner step with $K$ axpy operations. Without the bilinear structure the
naive evaluation of $\sum_jp_{ij}F(Y_j,Z)$ costs $i-1$ full evaluations
per inner step. For $K=0$ the combination is simply a weighted average
of the frozen fields, e.g.\
$\bar\theta_i=\sum_j\beta_{ij}\theta(Y_j)$, which is how split-explicit
codes are usually written.

\begin{remark}[consequences for the order conditions]
Since $\partial_z^2F=0$, all elementary differentials with a vertex that
has two or more white children vanish. Table~\ref{tab:treecount} lists
the number of remaining trees: at order 4 there are 20 instead of 26, at
order 5 there are 73 instead of 107. A method constructed for
\eqref{eq:B} only has to satisfy these conditions. Note that
\eqref{eq:B} contains \eqref{eq:L} and the linearized frame
representation \eqref{eq:F}.
\end{remark}

% =====================================================================
\section{MIS methods with polynomial coupling}
\label{sec:methods}

Throughout, $\tau\in[0,1]$ denotes the normalized time of the inner
problems.

\subsection{Definition}

\begin{definition}[MIS method with polynomial coupling]\label{def:mis}
Let $s\ge1$ and for $2\le i\le s+1$, $1\le j<i$ let $\alpha_{ij}\in\R$
and polynomials
\begin{equation*}
 p_{ij}(\tau)=\sum_{k=0}^{K}a_{ijk}\,\tau^k,\qquad
 g_{ij}(\tau)=\sum_{k=0}^{K}\gamma_{ijk}\,\tau^k
\end{equation*}
be given. One step $y_n\mapsto y_{n+1}$ is defined by $Y_1=y_n$ and for
$i=2,\dots,s+1$
\begin{align}
 Z_i(0) &= y_n+\sum_{j=1}^{i-1}\alpha_{ij}\,(Y_j-y_n), \label{eq:mis1}\\
 Z_i'(\tau) &= \sum_{j=1}^{i-1}g_{ij}(\tau)\,(Y_j-y_n)
   + h\sum_{j=1}^{i-1}p_{ij}(\tau)\,F\big(Y_j,Z_i(\tau)\big),
   \qquad \tau\in[0,1], \label{eq:mis2}\\
 Y_i &= Z_i(1), \label{eq:mis3}
\end{align}
and $y_{n+1}=Y_{s+1}$. We call $K$ the \emph{coupling degree}; $K=0$ is
\emph{constant coupling}. Constant $g_{ij}\equiv\gamma_{ij}$ gives the
generalized split-explicit methods of \cite{KnothWensch2014}.
\end{definition}

In the additive case \eqref{eq:A} the fast part is integrated with the
``speed'' $q_i(\tau):=\sum_jp_{ij}(\tau)$, i.e.\ over the fast time
$Q_i(1)$, $Q_i(\tau):=\int_0^\tau q_i$. We call
\begin{equation}\label{eq:internalconsistency}
 q_i(\tau)\equiv d_i \ \text{(constant)},\qquad\text{i.e.}\quad
 \sum_j a_{ij0}=d_i,\quad \sum_ja_{ijk}=0\ (k\ge1),
\end{equation}
\emph{uniform fast speed}. It is satisfied by all classical methods
below. Note that uniform fast speed does not make the inner problem
autonomous: the frozen coefficients $W_i(\tau)$ of
Lemma~\ref{lem:summation} still depend on $\tau$.

\begin{remark}
The coefficients used in \texttt{BSeries.jl}
(\texttt{MultirateInfinitesimalSplitMethod(A, D, G, c)}) correspond to
$A_{ij}=p_{ij}$ (polynomial matrix of size $(s+1)\times s$),
$D_{ij}=\alpha_{ij}$, and constant $G_{ij}=\gamma_{ij}$; the
implementation \cite{KetchesonRanocha2023} follows the algorithm of
Section~\ref{sec:algorithm}. Polynomial $g_{ij}$ require a one-line
change there (multiply by the polynomial before integrating).
\todo{cross-check conventions against the Fortran code
\texttt{MISBI\_Mod.f90}.}
\end{remark}

\subsection{Stage times and the base method}

Let $P_{ij}:=\int_0^1p_{ij}$ and $\Gamma_{ij}:=\int_0^1g_{ij}$. The
\emph{stage times} $c_i$ are defined by the method applied to $y'=1$:
\begin{equation}\label{eq:ci}
 c_1=0,\qquad c_i=\sum_{j<i}(\alpha_{ij}+\Gamma_{ij})\,c_j+\sum_{j<i}P_{ij}.
\end{equation}
If $F$ does not depend on its second argument, the inner problems are
pure quadratures and
$Y_i-y_n=\sum_j(\alpha_{ij}+\Gamma_{ij})(Y_j-y_n)+h\sum_jP_{ij}F(Y_j)$.
With the strictly lower triangular $(s+1)\times(s+1)$ matrices
$\bar A=(\alpha_{ij}+\Gamma_{ij})$ and $\bar P=(P_{ij})$ this is an
explicit Runge--Kutta method, the \emph{base method}, with
\begin{equation}\label{eq:basemethod}
 \begin{pmatrix}A^{RK}\\ b^{RK\,T}\end{pmatrix}=(I-\bar A)^{-1}\bar P .
\end{equation}
Consequently the conditions of all trees with only black vertices
(Section~\ref{sec:trees}) are exactly the order conditions of the base
method. The polynomial shape of $p_{ij}$ and $g_{ij}$ enters only the
trees with white vertices.

\subsection{Role of the polynomial coefficients}
\label{sec:polygamma}

Both kinds of coupling enter the inner problem as polynomials in $\tau$.
\begin{itemize}
\item \emph{$p_{ij}(\tau)$: frozen slow information.} The inner problem
  sees a time-dependent combination of the frozen right-hand sides, e.g.\
  a linear interpolation in time between stages (Example~\ref{ex:etd}).
  With $K=0$ the slow tendency is constant over the inner integration,
  as in classical split-explicit schemes.
\item \emph{$g_{ij}(\tau)$: increments of previous stages.} $Y_j-y_n$
  contains the fast evolution over $[0,c_jh]$, not only frozen slow
  information.
\item \emph{Separation of roles.} The shape of $p_{ij}$ and $g_{ij}$
  enters only the trees with white vertices; the base method sees only
  $\int_0^1p_{ij}$ and $\int_0^1g_{ij}$.
\item \emph{Cost.} Both terms are free in the implementation. Together
  they form a forcing $r_i(\tau)=\sum_{k=0}^K\tau^kR_{ik}$ with
  $R_{ik}=\sum_j\big(a_{ijk}f(Y_j)+\tfrac1h\gamma_{ijk}(Y_j-y_n)\big)$,
  $K+1$ vectors per stage (Section~\ref{sec:inner}).
\item \emph{Caution: $g_{ij}\ne0$ and stiffness.} For a stiff fast part,
  $Y_j-y_n$ is not $O(h)$. For $y'=Ly$ it equals $(e^{c_jhL}-I)y_n$,
  which is bounded but $O(1)$ for fast modes. The $g_{ij}$-terms then
  enter the stability function, generate extra stiff order conditions,
  and destroy the automatic exactness for the pure fast problem
  (Theorem~\ref{thm:exactfast}).
\end{itemize}
A natural strategy is therefore to construct first methods with
$g_{ij}\equiv0$ and $\alpha$ of the ``restart'' type of
Theorem~\ref{thm:exactfast}. The polynomial $g_{ij}$ (and general
$\alpha$) then serve as additional free parameters for stability
optimization, as the constant $\gamma_{ij}$ did in \cite{KnothWensch2014}.

\subsection{Classical methods as special cases}
\label{sec:specialcases}

\paragraph{Classical MIS methods.}
In \cite{WenschKnothGalant2009} the inner problem is posed on
$[0,d_ih]$ with constant coefficients $\beta_{ij}$ and
$d_i=\sum_j\beta_{ij}$. After rescaling the inner time to
$\tau\in[0,1]$ this is Definition~\ref{def:mis} with $K=0$,
$p_{ij}\equiv\beta_{ij}$, uniform fast speed $d_i$. The generalized
split-explicit methods of \cite{KnothWensch2014} use the additional
$\gamma_{ij}$ terms.

\paragraph{Split-explicit Runge--Kutta (Wicker--Skamarock).}
The RK3 scheme of \cite{WickerSkamarock2002} integrates every stage from
$y_n$ with the slow tendency of the previous stage held fixed:
$\alpha=0$, $g_{ij}\equiv0$, $c=(0,\tfrac13,\tfrac12,1)$,
$p_{i,i-1}\equiv c_i$, all other $p_{ij}=0$. It is of order 3 only for
linear problems; this is visible in the order conditions (the base method
is a three-stage RK with only a subdiagonal).

\paragraph{Telescopic (MRI-GARK type) methods.}
If $\alpha_{i,i-1}=1$ and all other $\alpha_{ij}=0$, $g=0$, then
$Z_i(0)=Y_{i-1}$ and the stages form a chain of inner integrations over
consecutive fast time intervals of length $d_i=c_i-c_{i-1}$. With
polynomial $p_{ij}$ these are the explicit multirate infinitesimal GARK
methods of Sandu \cite{Sandu2019} (in the additive setting); the uniform
fast speed \eqref{eq:internalconsistency} corresponds to its internal
consistency conditions. \todo{check correspondence of the coefficients
$\gamma_{ij}^{\{k\}}$ in \cite{Sandu2019} with $a_{ijk}$, including the
$1/k!$ or $1/h^k$ normalisation.} Related classes are the multirate
exponential RK (MERK) methods \cite{LuanChinomonaReynolds2020} and the
IMEX-MRI-GARK methods \cite{ChinomonaReynolds2021}.

\paragraph{Exponential Runge--Kutta methods.}
See Section~\ref{sec:exponential}: for the linear fast part \eqref{eq:L}
the inner problems can be solved by variation of constants, and MIS
methods with $\alpha=0$, $g_{ij}\equiv0$ and uniform fast speed $d_i=c_i$ are
exactly the exponential RK methods whose coefficients are combinations of
$\varphi_k(c_ihL)$.

\paragraph{Commutator-free Lie group methods.}
For the frame representation \eqref{eq:F}, $F(Y_j,\cdot)$ is the vector
field with coefficients frozen at $Y_j$. The exact flow over the time
$h\beta$ of the frozen field $\sum_j\beta_jF(Y_j,\cdot)$ is the
exponential used in commutator-free Lie group methods
\cite{CelledoniMarthinsenOwren2003,Owren2006}. A commutator-free method
that composes several such exponentials per stage is a telescopic MIS
method with constant coupling, where one stage of the Lie group method
corresponds to several MIS stages. With polynomial coupling one gets
``commutator-free methods with time-dependent frozen fields''.
\todo{work out the exact correspondence; do Owren's order conditions
coincide with ours restricted to \eqref{eq:F}?}

\paragraph{Semi-implicit / IMEX schemes.}
Schemes of the form \cite{BoscarinoFilbetRusso2016}
\begin{equation*}
 Y_i=y_n+\sum_{j<i}\hat a_{ij}k_j,\quad \tilde Z_i=y_n+\sum_{j<i}a_{ij}k_j,\quad
 k_i=hF\big(Y_i,\tilde Z_i+a_{ii}k_i\big),\quad y_{n+1}=y_n+\sum_ib_ik_i
\end{equation*}
also freeze the first argument and treat the second implicitly. They can
be read as MIS-type methods in which the inner problem is replaced by a
single implicit (DIRK) step: a case of \emph{non-exact} inner solution
(Section~\ref{sec:nonexact}).

\paragraph{Methods not contained.}
Integrating-factor (Lawson) methods \cite{Lawson1967} use
$e^{(c_i-c_j)hL}f(Y_j)$, i.e.\ the slow tendency is transported by the
fast flow. In MIS form this would require an impulsive forcing at
$\tau=c_j/c_i$, which is not polynomial. Such terms can at most be
generated through the initial value mixing $\alpha_{ij}$.

\begin{question}
Characterize the exponential integrators (Lawson, exponential RK with
mixed arguments $\varphi_k(c_jhL)$ in stage $i$, exponential Rosenbrock
methods \cite{HochbruckOstermannSchweitzer2009} with $L=\partial_y\mathcal
F(y_n)$) that are MIS methods in the sense of Definition~\ref{def:mis},
possibly with $\alpha\neq0$.
\end{question}

% =====================================================================
% =====================================================================
\section{Bicoloured trees and B-series}
\label{sec:trees}

\subsection{Trees and elementary differentials}

Let $\T$ be the set of rooted trees whose \emph{non-root} vertices are
coloured black (\,\Tb\,) or white. The colour of a vertex tells which
argument of the $F$ at the \emph{parent} vertex is differentiated:
black means the first (slow, frozen) argument, white the second (fast)
argument. Every tree is $t=[t_1,\dots,t_m]$, where the subtrees $t_k$
carry the colour $\kappa(t_k)\in\{\mathrm b,\mathrm w\}$ of their roots.
The \emph{elementary differentials} for \eqref{eq:G} are
\begin{equation*}
 \mathbf F(\bullet)=F(y,y),\qquad
 \mathbf F([t_1,\dots,t_m])
 =\partial_{\kappa(t_1)}\cdots\partial_{\kappa(t_m)}F(y,y)
  \big[\mathbf F(t_1),\dots,\mathbf F(t_m)\big],
\end{equation*}
with $\partial_{\mathrm b}=\partial_y$, $\partial_{\mathrm w}=\partial_z$.
Order $|t|$, symmetry $\sigma(t)$ (automorphisms preserving colours) and
density $\gamma(t)=|t|\prod_k\gamma(t_k)$ are defined as usual. Since
$y'=F(y,y)$, the exact solution is
\begin{equation*}
 y(t_n+h)=y_n+\sum_{t\in\T}\frac{h^{|t|}}{\sigma(t)\gamma(t)}\mathbf F(t)(y_n),
\end{equation*}
independent of the colours. The numbers of trees of order $1,\dots,6$
are $1,2,7,26,107,458$ (half the number of fully bicoloured trees). For
order $\le3$:
\begin{center}
\begin{tabular}{ccccccccccc}
\toprule
$t$ & \Tb & \Tbb & \Tbw & \Tbbb & \Tbbw & \Tbww & \Tbbbl & \Tbbwl & \Tbwbl & \Tbwwl\\
\midrule
$\gamma(t)$ & 1 & 2 & 2 & 3 & 3 & 3 & 6 & 6 & 6 & 6\\
$\sigma(t)$ & 1 & 1 & 1 & 2 & 1 & 2 & 1 & 1 & 1 & 1\\
\bottomrule
\end{tabular}
\end{center}
(the root is drawn black; its colour is irrelevant).

For the special cases the set of relevant trees shrinks:
\begin{itemize}
\item \eqref{eq:B}: $\partial_z^2F=0$, so a vertex has at most one
  white child (any number of black children).
\item \eqref{eq:A}: $\partial_y\partial_zF=0$, so all children of a
  vertex have the same colour.
\item \eqref{eq:L}: in addition $\partial_z^2F=0$. A vertex has either
  only black children or exactly one white child, so white vertices form
  ``$L$-chains''. These are the trees of exponential integrators.
\item Trees with only black vertices give the base method
  \eqref{eq:basemethod}. Trees with only white non-root vertices give the
  pure fast problem $y'=g(y)$.
\end{itemize}

\begin{table}[ht]
\centering
\caption{Number of trees (non-root vertices bicoloured, root colour
irrelevant) per order for the problem classes of
Section~\ref{sec:motivation}, i.e.\ number of order conditions at
each order. Computed by direct enumeration.}
\label{tab:treecount}
\begin{tabular}{lrrrrrrr}
\toprule
order & 1 & 2 & 3 & 4 & 5 & 6 & 7\\\midrule
general \eqref{eq:G} & 1 & 2 & 7 & 26 & 107 & 458 & 2058\\
affine in $z$ \eqref{eq:B} & 1 & 2 & 6 & 20 & 73 & 278 & 1106\\
additive \eqref{eq:A} & 1 & 2 & 6 & 18 & 60 & 204 & 734\\
constant linear \eqref{eq:L} & 1 & 2 & 5 & 13 & 37 & 108 & 332\\
classical RK (black only) & 1 & 1 & 2 & 4 & 9 & 20 & 48\\
\bottomrule
\end{tabular}
\end{table}

\subsection{B-series with polynomial coefficients}

Write the stage values and the inner solutions as B-series
\begin{equation*}
 Y_i=B(\eta_i,y_n),\qquad Z_i(\tau)=B\big(\zeta_i(\tau),y_n\big),
 \qquad B(a,y)=a(\emptyset)y+\sum_{t\in\T}\frac{h^{|t|}a(t)}{\sigma(t)}\mathbf F(t)(y),
\end{equation*}
with $\eta_i(\emptyset)=\zeta_i(\tau;\emptyset)=1$. The coefficients
$\zeta_i(\tau;t)$ are polynomials in $\tau$.

\begin{theorem}\label{thm:recursion}
The coefficients of the method of Definition~\ref{def:mis} satisfy, for
$t=[t_1,\dots,t_m]$, $2\le i\le s+1$, $\eta_1(t)=0$,
\begin{align*}
 \zeta_i(0;t) &= \sum_{j<i}\alpha_{ij}\,\eta_j(t),\\
 \frac{\dd}{\dd\tau}\zeta_i(\tau;t)
  &= \sum_{j<i}g_{ij}(\tau)\,\eta_j(t)
   + \sum_{j<i}p_{ij}(\tau)\prod_{k=1}^m
     \begin{cases}\eta_j(t_k), & \kappa(t_k)=\mathrm b,\\
                  \zeta_i(\tau;t_k), & \kappa(t_k)=\mathrm w,\end{cases}\\
 \eta_i(t) &= \zeta_i(1;t).
\end{align*}
The method has order $p$ if and only if $\eta_{s+1}(t)=1/\gamma(t)$ for
all $t\in\T$ with $|t|\le p$ (for generic $F$).
\end{theorem}

\begin{proof}[Sketch]
Insert the B-series of $Y_j$ and $Z_i(\tau)$ into $hF(Y_j,Z_i(\tau))$
and expand in a Taylor series in both arguments. The coefficient of
$\mathbf F([t_1,\dots,t_m])$ is the product of the coefficients of the
subtrees, taken from the series of the argument that is differentiated.
This is the standard P-series composition \cite{HairerNorsettWanner1993}.
Integrate in $\tau$ and use \eqref{eq:mis1}, \eqref{eq:mis3}.
\end{proof}

Since $\zeta_i(\tau;t)$ is a polynomial, the recursion is exact
rational arithmetic on polynomials. The degree in $\tau$ grows at most
like $\deg\zeta_i(\cdot;t)\le |t|(K+1)$.

% =====================================================================
\section{Order conditions}
\label{sec:orderconditions}

\subsection{Notation and general form}

For the coefficients of Definition~\ref{def:mis} define the polynomials
\begin{equation*}
 P_{ij}(\tau)=\int_0^\tau p_{ij},\qquad G_{ij}(\tau)=\int_0^\tau g_{ij},\qquad
 q_i=\sum_{j<i}p_{ij},\qquad Q_i=\sum_{j<i}P_{ij},
\end{equation*}
the strictly lower triangular $(s+1)\times(s+1)$ matrices
$\bar A=\big(\alpha_{ij}+G_{ij}(1)\big)$ and $\bar P=\big(P_{ij}(1)\big)$,
and
\begin{equation*}
 \Lambda=(I-\bar A)^{-1},\qquad \hat b^T=e_{s+1}^T\Lambda,\qquad
 \mathfrak A=\Lambda\bar P .
\end{equation*}
The stage times are $c=\Lambda\big(Q_i(1)\big)_i$ (cf.\ \eqref{eq:ci}),
and the \emph{inner time} of stage $i$ is the polynomial
\begin{equation*}
 z_i(\tau):=\zeta_i(\tau;\bullet)=\sum_{j<i}\big(\alpha_{ij}+G_{ij}(\tau)\big)c_j+Q_i(\tau),
 \qquad z_i(0)=\sum_j\alpha_{ij}c_j,\quad z_i(1)=c_i .
\end{equation*}
$\mathfrak A$ is the coefficient matrix of the base method
\eqref{eq:basemethod}: its rows $2,\dots,s$ are the stage coefficients,
its last row is $b^{RK\,T}$.

\begin{theorem}[general form of the order conditions]\label{thm:general}
For $t=[t_1,\dots,t_m]$ let
\begin{equation}\label{eq:dit}
 d_i(t)=\sum_{j<i}\int_0^1p_{ij}(\sigma)
   \prod_{k:\,\kappa(t_k)=\mathrm b}\eta_j(t_k)
   \prod_{k:\,\kappa(t_k)=\mathrm w}\zeta_i(\sigma;t_k)\dd\sigma .
\end{equation}
Then
\begin{equation*}
 \eta(t)=\Lambda\,d(t),\qquad
 \zeta_i(\tau;t)=\sum_{j<i}\big(\alpha_{ij}+G_{ij}(\tau)\big)\eta_j(t)
   +\sum_{j<i}\int_0^\tau p_{ij}(\sigma)\prod\cdots\dd\sigma,
\end{equation*}
with the same products as in \eqref{eq:dit}. The method has order $p$ if
and only if
\begin{equation}\label{eq:oc}
 \hat b^Td(t)=\sum_{i=2}^{s+1}\hat b_i\,d_i(t)=\frac1{\gamma(t)}
 \qquad\text{for all }t\in\T,\ |t|\le p .
\end{equation}
\end{theorem}

\begin{proof}
Integrating the recursion of Theorem~\ref{thm:recursion} from $0$ to $1$
gives $\eta_i(t)=\sum_j\bar A_{ij}\eta_j(t)+d_i(t)$, i.e.\
$(I-\bar A)\eta(t)=d(t)$. The formula for $\zeta_i$ is the integral from
$0$ to $\tau$.
\end{proof}

The structure is that of Runge--Kutta conditions $b^T\Phi(t)=1/\gamma(t)$:
\begin{itemize}
\item $\hat b$ plays the role of the weights;
\item $d_i(t)$ is a ``stage increment'' that integrates the product of
  the subtree coefficients against the coupling polynomials;
\item black subtrees contribute stage values $\eta_j$ of the
  \emph{coupled} stage $j$;
\item white subtrees contribute the \emph{inner} coefficients
  $\zeta_i(\sigma;\cdot)$ of the \emph{current} stage $i$, evaluated
  along the inner time.
\end{itemize}
Since all quantities are polynomials, every integral is evaluated
exactly with $\int_0^1\sigma^k\dd\sigma=1/(k+1)$. Every condition is
therefore a polynomial equation in $\alpha_{ij}$, $a_{ijk}$,
$\gamma_{ijk}$.

\subsection{Explicit conditions up to order 3}

Table~\ref{tab:oc3} lists $d_i(t)$ for all ten trees of order $\le3$; the
conditions are $\hat b^Td(t)=1/\gamma(t)$. The auxiliary stage
quantities are
\begin{align*}
 \eta(\Tbb)&=\Lambda\Big(\sum_jP_{ij}(1)c_j\Big)_i=\mathfrak Ac, &
 \eta(\Tbw)&=\Lambda\Big(\int_0^1q_iz_i\Big)_i,\\
 \zeta_i(\sigma;\Tbb)&=\sum_j\big(\alpha_{ij}+G_{ij}(\sigma)\big)\eta_j(\Tbb)+\sum_jP_{ij}(\sigma)c_j, &
 \zeta_i(\sigma;\Tbw)&=\sum_j\big(\alpha_{ij}+G_{ij}(\sigma)\big)\eta_j(\Tbw)+\int_0^\sigma q_iz_i .
\end{align*}

\begin{table}[ht]
\centering
\caption{Order conditions $\sum_i\hat b_id_i(t)=1/\gamma(t)$ up to
order 3 for the general method of Definition~\ref{def:mis}. Trees are
drawn with the root at the bottom (root colour irrelevant). Column ``class'' refers
to Theorem~\ref{thm:exactfast}: B = base method condition, F =
automatically satisfied, C = coupling condition, -- = implied.}
\label{tab:oc3}
\renewcommand{\arraystretch}{1.6}
\begin{tabular}{rcllc}
\toprule
order & $t$ & $d_i(t)$ & $1/\gamma(t)$ & class\\\midrule
1 & \Tb & $Q_i(1)$ & $1$ & B\\\midrule
2 & \Tbb & $\sum_jP_{ij}(1)\,c_j$ & $\frac12$ & B\\
  & \Tbw & $\int_0^1q_i(\sigma)\,z_i(\sigma)\dd\sigma$ & $\frac12$ & F\\\midrule
3 & \Tbbb & $\sum_jP_{ij}(1)\,c_j^2$ & $\frac13$ & B\\
  & \Tbbw & $\sum_jc_j\int_0^1p_{ij}(\sigma)\,z_i(\sigma)\dd\sigma$ & $\frac13$ & C\\
  & \Tbww & $\int_0^1q_i(\sigma)\,z_i(\sigma)^2\dd\sigma$ & $\frac13$ & F\\
  & \Tbbbl & $\sum_jP_{ij}(1)\,\eta_j(\Tbb)$ & $\frac16$ & B\\
  & \Tbbwl & $\sum_jP_{ij}(1)\,\eta_j(\Tbw)$ & $\frac16$ & --\\
  & \Tbwbl & $\int_0^1q_i(\sigma)\,\zeta_i(\sigma;\Tbb)\dd\sigma$ & $\frac16$ & --\\
  & \Tbwwl & $\int_0^1q_i(\sigma)\,\zeta_i(\sigma;\Tbw)\dd\sigma$ & $\frac16$ & F\\
\bottomrule
\end{tabular}
\end{table}

The order 4 conditions (26 trees) are listed in
Appendix~\ref{app:order4}. They were generated from \eqref{eq:dit}. All
formulas of Table~\ref{tab:oc3} were checked against the recursion of
Theorem~\ref{thm:recursion} for random methods with general $\alpha$,
$g_{ij}$ and $K=2$ (Section~\ref{sec:algorithm}).

\subsection{Special variants and their order condition theorems}
\label{sec:specialtheorems}

\paragraph{Base method.}
If $F$ does not depend on $z$, only black trees remain and
\eqref{eq:oc} reduces to the Runge--Kutta conditions of $\mathfrak A$:

\begin{proposition}\label{prop:base}
For trees with only black vertices, $\eta(t)$ is the Runge--Kutta
elementary weight of the base method $\mathfrak A=\Lambda\bar P$. In
particular, order $p$ of the MIS method implies order $p$ of the base
method.
\end{proposition}

\paragraph{Methods that are exact for the fast part.}
The most important subclass consists of methods whose inner problems
start at a previous stage and contain no $(Y_j-y_n)$ terms:
\begin{equation}\label{eq:restart}
 g_{ij}\equiv0,\qquad Z_i(0)=Y_{k(i)}\ \text{for some }k(i)<i
 \quad(\text{i.e.\ row $i$ of $\alpha$ is $e_{k(i)}^T$, or $0$ if $k(i)=1$}).
\end{equation}
It contains:
\begin{itemize}
\item the classical MIS methods with $Z_i(0)=y_n$;
\item the telescopic methods with $Z_i(0)=Y_{i-1}$, in particular
  MRI-GARK \cite{Sandu2019};
\item the exponential RK methods of Section~\ref{sec:exponential};
\item arbitrary polynomial coupling $p_{ij}$ (uniform fast speed is
  \emph{not} required).
\end{itemize}
Under \eqref{eq:restart}:
\begin{itemize}
\item $\Lambda$ has a simple structure, $\eta_i=\eta_{k(i)}+d_i$;
\item hence $\hat b_i=1$ for the stages on the \emph{chain}
  $s+1\to k(s+1)\to k(k(s+1))\to\dots\to1$, and $\hat b_i=0$ otherwise;
\item the inner time is $z_i(\sigma)=c_{k(i)}+Q_i(\sigma)$.
\end{itemize}

\begin{theorem}\label{thm:exactfast}
Assume \eqref{eq:restart} and $c_{s+1}=1$.
\begin{enumerate}[label=(\alph*)]
\item For every tree $u$ whose non-root vertices are all white,
  $\eta_i(u)=c_i^{|u|}/\gamma(u)$ and
  $\zeta_i(\sigma;u)=z_i(\sigma)^{|u|}/\gamma(u)$. In particular all
  pure fast conditions (class F) are satisfied.
\item For all stages
  \begin{equation*}
   \eta_i(\Tbbwl)=\tfrac12\,\eta_i(\Tbbb),\qquad
   \eta_i(\Tbbw)+\eta_i(\Tbwbl)=c_i\,\eta_i(\Tbb).
  \end{equation*}
\item The method has order 3 if and only if the base method $\mathfrak A$
  has order 3 and the single coupling condition
  \begin{equation}\label{eq:coupling3}
   \sum_{i\in\text{chain}}\ \sum_{j<i}c_j\int_0^1p_{ij}(\sigma)\big(c_{k(i)}+Q_i(\sigma)\big)\dd\sigma=\frac13
  \end{equation}
  holds.
\end{enumerate}
\end{theorem}

\begin{proof}
(a) For $F(y,z)=g(z)$ the inner problem is $Z_i'=hq_i(\tau)g(Z_i)$,
$Z_i(0)=Y_{k(i)}$. This is the flow of $g$ reparametrized in time; hence
$Z_i(\sigma)$ is the exact flow over the time $h(c_{k(i)}+Q_i(\sigma))=hz_i(\sigma)$,
provided $Y_{k(i)}$ is the exact flow over $hc_{k(i)}$. Induction over
$i$ gives the claim for $\zeta$ and $\eta$. The coefficients of the
exact flow over time $hc$ are $c^{|u|}/\gamma(u)$.

(b) By (a), $\eta_j(\Tbw)=c_j^2/2$, so
$d_i(\Tbbwl)=\tfrac12d_i(\Tbbb)$. For the second identity, by (a) and
integration by parts with $Q_i'=q_i$, $P_{ij}'=p_{ij}$:
\begin{align*}
 d_i(\Tbbw)+d_i(\Tbwbl)
 &=\sum_jc_j\int_0^1p_{ij}\big(c_{k}+Q_i\big)
  +\int_0^1q_i\Big(\eta_{k}(\Tbb)+\sum_jP_{ij}c_j\Big)\\
 &=c_k\sum_jP_{ij}(1)c_j+\sum_jc_j\big[P_{ij}Q_i\big]_0^1+Q_i(1)\eta_k(\Tbb)\\
 &=c_i\big(\eta_i(\Tbb)-\eta_k(\Tbb)\big)+(c_i-c_k)\eta_k(\Tbb)
  =c_i\eta_i(\Tbb)-c_k\eta_k(\Tbb),
\end{align*}
with $k=k(i)$, $c_i=c_k+Q_i(1)$ and
$\eta_i(\Tbb)=\eta_k(\Tbb)+\sum_jP_{ij}(1)c_j$. Since
$\eta_i=\eta_k+d_i$, induction over $i$ gives the identity.

(c) By (a) the conditions for \Tbw, \Tbww, \Tbwwl\ hold. By (b),
$\eta_{s+1}(\Tbbwl)=\tfrac16\iff\eta_{s+1}(\Tbbb)=\tfrac13$, and, given
$\eta_{s+1}(\Tbb)=\tfrac12$,
$\eta_{s+1}(\Tbwbl)=\tfrac16\iff\eta_{s+1}(\Tbbw)=\tfrac13$, which is
\eqref{eq:coupling3}. The remaining conditions \Tb, \Tbb, \Tbbb, \Tbbbl\
are those of the base method.
\end{proof}

\begin{proposition}[computational]\label{prop:count}
Under \eqref{eq:restart}, the number of coupling conditions that are
independent of the base method conditions and of the lower order
conditions is
\begin{center}
\begin{tabular}{lcccc}
\toprule
order & 2 & 3 & 4 & 5\\\midrule
all trees & 2 & 7 & 26 & 107\\
base method (black trees) & 1 & 2 & 4 & 9\\
coupling conditions, class \eqref{eq:restart} & 0 & 1 & 4 & 16\\
coupling conditions, general $\alpha$, $g_{ij}$ & 1 & 5 & 22 & 89\\
\bottomrule
\end{tabular}
\end{center}
At order 4 a possible choice of the four coupling conditions is
(with $\hat b_i=1$ on the chain, $k=k(i)$)
\begin{align*}
 &\Tfourbbw: && \sum_{i}\hat b_i\sum_jc_j^2\int_0^1p_{ij}(\sigma)\,z_i(\sigma)\dd\sigma=\tfrac14, &
 &\Tfourbww: && \sum_{i}\hat b_i\sum_jc_j\int_0^1p_{ij}(\sigma)\,z_i(\sigma)^2\dd\sigma=\tfrac14,\\
 &\Tfourbbl: && \sum_{i}\hat b_i\sum_j\eta_j(\Tbb)\int_0^1p_{ij}(\sigma)\,z_i(\sigma)\dd\sigma=\tfrac18, &
 &\Tfourbbwl: && \sum_{i}\hat b_i\sum_jP_{ij}(1)\,\eta_j(\Tbbw)=\tfrac1{12},
\end{align*}
with $z_i(\sigma)=c_{k}+Q_i(\sigma)$ and
$\eta_j(\Tbbw)=\sum_{l\in\text{chain}(j)}\sum_rc_r\int_0^1p_{lr}z_l$.
All other mixed trees of order 4 are implied (class -- in
Table~\ref{tab:oc4}).
\end{proposition}

The counts were obtained as ranks of the stage-wise coefficient vectors
of many random methods. The ranks were computed exactly modulo a large
prime, with the products of lower-order coefficients included as
admissible relations. For orders 2 and 3 they agree with
Theorem~\ref{thm:exactfast}. A proof of the order 4 statement in the
style of Theorem~\ref{thm:exactfast}(b) is open.

\paragraph{Remarks on known results.}
\begin{itemize}
\item \emph{MRI-GARK.} Sandu \cite{Sandu2019} derived order conditions
  for MRI-GARK methods (telescopic, polynomial coupling, additive
  splitting) via GARK theory: the conditions of the slow base method plus
  coupling conditions, under ``internal consistency'' (uniform fast
  speed). Theorem~\ref{thm:exactfast} shows that for exact inner
  solutions uniform fast speed is not needed for the order. It remains
  relevant for the inner integrator, since $q_i<0$ would mean integrating
  the fast part backwards in time. As a check, our implementation gives
  order 3 for MRI-GARK-ERK33 and order 4 for MRI-GARK-ERK45a (coefficients
  from \texttt{MISBI\_Mod.f90}). \todo{compare the number of coupling
  conditions with \cite{Sandu2019}.}
\item \emph{Classical MIS and generalized split-explicit methods.} For
  constant coupling the conditions up to order 3 are given in
  \cite{WenschKnothGalant2009,KnothWensch2014} and extended in
  \cite{BauerKnoth2021}; Table~\ref{tab:oc3} reduces to them for $K=0$.
  \todo{check notation correspondence.}
\item \emph{Additive splitting \eqref{eq:A}.} Only trees in which all
  children of a vertex have the same colour occur
  (Table~\ref{tab:treecount}). These are the GARK conditions
  \cite{GuntherSandu2016} of the equivalent additive method.
\item \emph{Constant linear fast part \eqref{eq:L}.} Only trees whose
  white vertices form chains occur. These are the classical (non-stiff)
  conditions of exponential RK methods; the stiff conditions are treated
  in Section~\ref{sec:exponential}.
\item \emph{Frame representation \eqref{eq:F}.} The conditions for
  commutator-free Lie group methods \cite{Owren2006} should arise from
  \eqref{eq:oc} restricted to the trees of \eqref{eq:F}.
  \todo{verify.}
\item \emph{Affine fast part \eqref{eq:B}.} Trees with a vertex with two
  or more white children drop out (20 instead of 26 at order 4).
\end{itemize}

\begin{question}
Is there a simplifying assumption, analogous to the stage order
condition $C(q)$ of RK methods, that reduces the conditions of
Theorem~\ref{thm:recursion} to those of the base method plus a small
number of coupling conditions? A natural candidate is the polynomial
identity
\begin{equation}\label{eq:Cq}
 \sum_{j<i}p_{ij}(\tau)\,c_j^{k-1}=c_i^{k}\,\tau^{k-1},\qquad k=1,\dots,q,
\end{equation}
(for $\alpha=0$, $g_{ij}\equiv0$), which appears again as the strong form of the
stiff order conditions in Section~\ref{sec:exponential}. For explicit
methods \eqref{eq:Cq} cannot hold for $i=2$ and $k\ge2$, as for RK
methods.
\end{question}

\section{Computing the order conditions}
\label{sec:algorithm}

\subsection{Algorithm}

Theorem~\ref{thm:recursion} gives an algorithm that is independent of the
special structure and works with exact rational (or symbolic)
arithmetic. The trees are processed in increasing order, so all subtree
coefficients are available.

\begin{center}
\begin{minipage}{0.9\linewidth}
\begin{tabbing}
xx\=xx\=xx\=xx\=xx\=\kill
\textbf{for} each tree $t=[t_1,\dots,t_m]$, ordered by $|t|$\\
\>\textbf{for} $i=2,\dots,s+1$\\
\>\>$\phi(\tau)\gets\sum_{j<i}g_{ij}(\tau)\eta_j(t)$\\
\>\>\textbf{for} $j=1,\dots,i-1$\\
\>\>\>$r(\tau)\gets1$\\
\>\>\>\textbf{for} $k=1,\dots,m$: \ $r\gets r\cdot\zeta_i(\cdot;t_k)$ if
       $t_k$ white (polynomial product), else $r\gets r\cdot\eta_j(t_k)$\\
\>\>\>$\phi\gets\phi+p_{ij}\,r$\\
\>\>$\zeta_i(\tau;t)\gets\sum_{j<i}\alpha_{ij}\eta_j(t)+\int_0^\tau\phi(\tilde\tau)\dd\tilde\tau$\\
\>\>$\eta_i(t)\gets\zeta_i(1;t)$\\
\textbf{residual:} $\eta_{s+1}(t)-1/\gamma(t)$
\end{tabbing}
\end{minipage}
\end{center}

Remarks:
\begin{itemize}
\item This is implemented in the Fortran program (repository
  \texttt{OsKnoth/B-Series}) up to order 4 and in \texttt{bseries} for
  \texttt{MultirateInfinitesimalSplitMethod} in \texttt{BSeries.jl}
  \cite{KetchesonRanocha2023} for arbitrary order.
\item With symbolic coefficients (SymPy, SymEngine, Symbolics via the
  extensions of \texttt{BSeries.jl}) the residuals are polynomial
  equations in $a_{ijk},\alpha_{ij},\gamma_{ijk}$. These can be printed as
  \LaTeX\ tables and are the starting point for constructing new methods.
\item The algorithm needs only the structure ``which argument is
  differentiated''. Restricting to the trees of \eqref{eq:A} or
  \eqref{eq:L} gives the conditions for those classes.
\end{itemize}

\subsection{Example (sketch, \texttt{BSeries.jl})}

{\small
\begin{verbatim}
using BSeries, RootedTrees, Polynomials
P(v) = Polynomial{Rational{Int}, :τ}(v)
# ETD2RK (Cox--Matthews) as MIS method with linear-in-time coupling
A = [P([0])     P([0]);
     P([1])     P([0]);
     P([1, -1]) P([0, 1])]          # p_ij(τ), size (s+1) x s
D = zeros(Rational{Int}, 3, 2)      # alpha
G = zeros(Rational{Int}, 3, 2)      # gamma
c = [0, 1]
mis = MultirateInfinitesimalSplitMethod(A, D, G, c)
series = bseries(mis, 4)
for (t, v) in series
    RootedTrees.order(t) == 0 && continue
    res = v - 1 // RootedTrees.density(t)
    iszero(res) || println(t, "  residual = ", res)
end
\end{verbatim}
}
\noindent The same computation with the Python reference implementation
(see below) gives order 2. The residuals
$\eta_{s+1}(t)-1/\gamma(t)$ of order 3 are
\begin{center}
\begin{tabular}{cccccccc}
\toprule
$t$ & \Tbbb & \Tbbw & \Tbww & \Tbbbl & \Tbbwl & \Tbwbl & \Tbwwl\\\midrule
residual & $\frac16$ & $0$ & $0$ & $-\frac16$ & $\frac1{12}$ & $0$ & $0$\\
\bottomrule
\end{tabular}
\end{center}
In accordance with Theorem~\ref{thm:exactfast}:
\begin{itemize}
\item the pure fast conditions and the coupling condition
  \eqref{eq:coupling3} hold;
\item the order is limited by the base method, which is Heun's method
  ($b=(\frac12,\frac12)$, $a_{21}=1$);
\item the residual of \Tbbwl\ is half that of \Tbbb.
\end{itemize}

\paragraph{Python reference implementation.}
For checking the formulas of this paper we wrote a short, independent
implementation of Theorem~\ref{thm:recursion} in exact rational
arithmetic (\texttt{misbseries.py}). It covers general $\alpha$,
polynomial $g_{ij}$ and $p_{ij}$, and arbitrary order. It reproduces:
\begin{itemize}
\item order 3 for MRI-GARK-ERK33 and order 4 for MRI-GARK-ERK45a;
\item Tables~\ref{tab:oc3} and \ref{tab:oc4};
\item Theorem~\ref{thm:exactfast} (checked up to order 5);
\item the counts of Proposition~\ref{prop:count} and
  Table~\ref{tab:treecount}.
\end{itemize}

\subsection{Planned computations (order conditions only)}

\begin{enumerate}
\item Cross-check of the Fortran program and \texttt{BSeries.jl} for all
  methods in \texttt{MISBI\_Mod.f90} up to order 4; extension to order 5.
\item Tests of Proposition~\ref{prop:exprk} (exponential RK methods from
  the literature); proof of Proposition~\ref{prop:count} at order 4.
\item Symbolic conditions for small families (e.g.\ $s=2,3$, $K\le2$)
  for \eqref{eq:G}, \eqref{eq:A}, \eqref{eq:L}. Count of independent
  conditions per order and class.
\item The order conditions with non-exact inner solution
  (Section~\ref{sec:nonexact}).
\item The conditions restricted to the affine class \eqref{eq:B}
  (Table~\ref{tab:treecount}); polynomial $g_{ij}$ in the recursion.
\end{enumerate}

\section{Linear fast part: exponential integrators and stiff order conditions}
\label{sec:exponential}

\subsection{Closed form of the stages}

For \eqref{eq:L} with $\alpha=0$, $g_{ij}\equiv0$ and uniform fast speed $q_i\equiv
c_i$ the inner problem
$Z_i'=h\sum_jp_{ij}(\tau)f(Y_j)+hc_iLZ_i$, $Z_i(0)=y_n$, is solved by
variation of constants. With
\begin{equation*}
 \varphi_{k+1}(z)=\int_0^1e^{(1-\tau)z}\frac{\tau^k}{k!}\dd\tau,\qquad k\ge0,
\end{equation*}
one obtains
\begin{equation}\label{eq:exprk}
 Y_i=e^{c_ihL}y_n+h\sum_{j<i}a_{ij}(hL)f(Y_j),\qquad
 a_{ij}(z)=\int_0^1e^{(1-\tau)c_iz}p_{ij}(\tau)\dd\tau
 =\sum_{k=0}^{K_i}k!\,a_{ijk}\,\varphi_{k+1}(c_iz).
\end{equation}

\begin{proposition}\label{prop:exprk}
An exponential RK method
$Y_i=e^{c_ihL}y_n+h\sum_ja_{ij}(hL)f(Y_j)$, $y_{n+1}=Y_{s+1}$, $c_{s+1}=1$,
is an MIS method with polynomial coupling ($\alpha=0$, $g_{ij}\equiv0$) if and only
if every $a_{ij}(z)$ is a linear combination of $\varphi_k(c_iz)$,
$k\ge1$, and $\sum_ja_{ij}(z)=c_i\varphi_1(c_iz)$.
\end{proposition}

\begin{proof}
By \eqref{eq:exprk}, using the linear independence of
$\varphi_1(cz),\varphi_2(cz),\dots$ for fixed $c\neq0$. The condition
$\sum_ja_{ij}=c_i\varphi_1(c_iz)$ is equivalent to
$q_i(\tau)=\sum_jp_{ij}(\tau)\equiv c_i$.
\end{proof}

Hence such methods, applied to \eqref{eq:G} instead of \eqref{eq:L},
are a natural nonlinear generalization of exponential RK methods, in
which the exponential is replaced by the flow of the fast part with
frozen slow argument.

\begin{example}\label{ex:etd}
\begin{enumerate}[label=(\alph*)]
\item Exponential Euler, $y_{n+1}=e^{hL}y_n+h\varphi_1(hL)f(y_n)$:
  $s=1$, $p_{21}=1$. In MIS form this is the ``frozen tendency''
  step.
\item ETD2RK of Cox and Matthews \cite{CoxMatthews2002}: $c_2=1$,
  $a_{21}=\varphi_1$, $b_1=\varphi_1-\varphi_2$, $b_2=\varphi_2$. In MIS
  form
  \begin{equation*}
   p_{21}=1,\qquad p_{31}(\tau)=1-\tau,\qquad p_{32}(\tau)=\tau:
  \end{equation*}
  the slow tendency is \emph{linearly interpolated in time} between the
  two stages, which is the simplest genuinely polynomial MIS method.
\item The two-stage family of \cite{HochbruckOstermann2005} with
  $a_{21}=c_2\varphi_1(c_2z)$, $b_1=\varphi_1-\varphi_2/c_2$,
  $b_2=\varphi_2/c_2$: $p_{21}=c_2$, $p_{31}=1-\tau/c_2$,
  $p_{32}=\tau/c_2$.
\end{enumerate}
\end{example}

Without uniform fast speed the closed form remains available:
\begin{equation*}
Y_i=e^{Q_i(1)hL}y_n+h\sum_j\int_0^1e^{(Q_i(1)-Q_i(\tau))hL}p_{ij}(\tau)\dd\tau\,f(Y_j).
\end{equation*}
The weights are then no longer combinations of $\varphi$-functions. With
$\alpha\neq0$ (in particular telescopic methods) products
$e^{d_iz}\varphi_k(c_jz)$ appear, i.e.\ exponential RK methods with mixed
arguments.

\subsection{Stiff order conditions}

The classical order conditions assume that $h\|\partial_zF\|$ is small.
In NWP this is not the case: $hL$ has large eigenvalues (on or near the
imaginary axis for acoustic waves). For exponential RK methods,
Hochbruck and Ostermann \cite{HochbruckOstermann2005,HochbruckOstermann2010}
derived \emph{stiff order conditions}, which must hold uniformly in $hL$.
With
\begin{equation*}
 \psi_j(z)=\sum_ib_i(z)\frac{c_i^{j-1}}{(j-1)!}-\varphi_j(z),\qquad
 \psi_{j,i}(z)=\sum_ka_{ik}(z)\frac{c_k^{j-1}}{(j-1)!}-c_i^j\varphi_j(c_iz),
\end{equation*}
($b_i=a_{s+1,i}$) the conditions up to order 4 read
\begin{center}
\begin{tabular}{cl}
\toprule
order & condition\\\midrule
1 & $\psi_1(hL)=0$\\
2 & $\psi_2(hL)=0$, \quad $\psi_{1,i}(hL)=0$\\
3 & $\psi_3(hL)=0$, \quad $\sum_ib_i(hL)J\psi_{2,i}(hL)=0$\\
4 & $\psi_4(hL)=0$, \quad $\sum_ib_i(hL)J\psi_{3,i}(hL)=0$,\\
  & $\sum_ib_i(hL)J\sum_ka_{ik}(hL)J\psi_{2,k}(hL)=0$, \quad
    $\sum_ib_i(hL)c_iK\psi_{2,i}(hL)=0$\\
\bottomrule
\end{tabular}
\end{center}
where $J,K$ are arbitrary bounded operators. \todo{check table against
\cite{HochbruckOstermann2005}.}

For MIS methods with polynomial coupling all these conditions become
\emph{algebraic} conditions on the coefficients $a_{ijk}$:

\begin{proposition}\label{prop:stiff}
Let $\alpha=0$, $g_{ij}\equiv0$, $q_i\equiv c_i$, and define the polynomials
\begin{equation*}
 r_i^{[j]}(\tau)=\sum_{k<i}p_{ik}(\tau)\frac{c_k^{j-1}}{(j-1)!}
   -c_i^j\frac{\tau^{j-1}}{(j-1)!} .
\end{equation*}
\begin{enumerate}[label=(\alph*)]
\item $\psi_{j,i}\equiv0 \iff r_i^{[j]}\equiv0$, and
  $\psi_j\equiv0\iff r_{s+1}^{[j]}\equiv0$. In particular $\psi_{1,i}=0$
  is the uniform fast speed condition.
\item A sufficient condition for $\sum_ib_i(z)J\psi_{2,i}(z)=0$ for all
  $J$ is that for every value $c$ of the stage times
  \begin{equation*}
   \sum_{i:\,c_i=c}p_{s+1,i}(\tau)\,r_i^{[2]}(\sigma)\equiv0
   \qquad\text{as a polynomial in }(\tau,\sigma).
  \end{equation*}
  Analogous bivariate (resp.\ trivariate, grouped by pairs
  $(c_i,c_k)$) polynomial identities are sufficient for the order 4
  conditions.
\end{enumerate}
\end{proposition}

\begin{proof}[Sketch]
(a) $\psi_{j,i}(z)=\int_0^1e^{(1-\tau)c_iz}r_i^{[j]}(\tau)\dd\tau$,
and the map from polynomials to these functions is injective for
$c_i\neq0$. (b) $b_i(z_1)\psi_{2,i}(z_2)
=\iint e^{(1-\tau)z_1}e^{(1-\sigma)c_iz_2}p_{s+1,i}(\tau)r_i^{[2]}(\sigma)\dd\tau\dd\sigma$;
since $J$ is arbitrary, the two arguments must be treated as independent.
\end{proof}

Example: for ETD2RK, $r_3^{[2]}=(1-\tau)\cdot0+\tau\cdot1-\tau=0$,
so it has stiff order 2.

\begin{question}
Are the sufficient conditions of Proposition~\ref{prop:stiff}(b) also
necessary, i.e.\ are the functions $\varphi_k(cz)$ for different $c$ and
$k$ linearly independent in the required sense? How many stages and which
coupling degree are needed for stiff order 3 and 4? (For exponential RK,
stiff order 4 needs five stages \cite{HochbruckOstermann2005}.)
\end{question}

\begin{question}
The theory of \cite{HochbruckOstermann2005} is formulated for parabolic
problems (analytic semigroups), where parabolic smoothing allows weakened
conditions. Acoustic operators are skew-adjoint (or nearly so), with
$\|e^{tL}\|\le1$ but no smoothing. Which conditions are needed in the
hyperbolic case? Is there an analogue of stiff order conditions for the
nonlinear fast part \eqref{eq:G}, e.g.\ with B-series whose coefficients
are functions of $h\partial_zF$ (resummation of the $L$-chains)?
\end{question}

% =====================================================================
\section{Spatial discretizations}
\label{sec:discretization}

The freezing of Section~\ref{sec:frozen} is applied to the
semi-discrete equations. We discuss two standard discretization types.

\subsection{Finite volumes with high-order advection}

Typical NWP dynamical cores (WRF, COSMO, ICON-like) use a staggered
Arakawa C-grid with face-normal momenta. In one direction, with cell
index $i$ and face mass fluxes $M_{i+1/2}=(\rho u)_{i+1/2}$:
\begin{align*}
 \dot\rho_i&=-\tfrac1{\Delta x}\big(M_{i+1/2}-M_{i-1/2}\big),\\
 \dot{(\rho\theta)}_i&=-\tfrac1{\Delta x}\big(M_{i+1/2}\,\hat\theta_{i+1/2}-M_{i-1/2}\,\hat\theta_{i-1/2}\big),\\
 \dot M_{i+1/2}&=\mathrm{Adv}_{i+1/2}(y)-\tfrac1{\Delta x}\big(p_{i+1}-p_i\big)+\dots
\end{align*}
\begin{itemize}
\item The face value $\hat\theta_{i+1/2}=\mathcal R(\theta;\operatorname{sign}u_{i+1/2})$
  is a high-order upwind-biased reconstruction: 3rd/5th-order linear
  upwind \cite{WickerSkamarock2002} or WENO \cite{JiangShu1996}. It is
  the expensive, nonlinear part.
\item In the frozen form $\hat\theta$ (including the upwind direction and
  the nonlinear WENO weights) is computed once per stage from $Y_j$; the
  mass fluxes $M$ are fast. The inner $\rho\theta$ equation is then a
  first-order flux-form operator with prescribed face values, and
  $\sum_jp_{ij}(\tau)\hat\theta_j$ is a face field.
\item Option A gives
  $-a_{i+1/2}(y)\big((\rho\theta)_{i+1}-(\rho\theta)_i\big)/\Delta x$
  with $a$ averaged to faces. Option B gives
  $-\big(\pi_{i+1}(\rho\theta)_{i+1}-\pi_i(\rho\theta)_i\big)/\Delta x$
  with cell values $\pi$.
\item In the vertical, the acoustic and buoyancy terms are usually
  treated implicitly in the inner problem (HEVI); each column is then a
  tridiagonal solve.
\item Terrain-following coordinates add metric terms, which are linear in
  the fast variables and do not change the structure.
\end{itemize}

A useful property: if $\theta\equiv\theta_0$ is constant, then
$\hat\theta_j\equiv\theta_0$ for all $j$ and
$\sum_jp_{ij}(\tau)\hat\theta_j=q_i(\tau)\theta_0$. The inner $\rho\theta$
equation is then $\theta_0$ times the inner mass equation, so constant
$\theta$ is preserved exactly (tracer--mass consistency) for every
method and every inner integrator that treats both equations alike.

\subsection{Discontinuous Galerkin methods}

Nodal DG methods on Legendre--Gauss--Lobatto nodes with the
summation-by-parts property and split-form (flux differencing) volume
terms \cite{GassnerWintersKopriva2016} are increasingly used in
atmospheric dynamical cores \cite{GiraldoKellyConstantinescu2013}. In
one element and one direction:
\begin{equation*}
 \dot{\mathbf q}_l=-\sum_m2D_{lm}\,\mathbf f^\#(\mathbf q_l,\mathbf q_m)
 -\frac{\delta_{l\in\partial}}{\omega_l}\big(\mathbf f^*(\mathbf q^-,\mathbf q^+)-\mathbf f(\mathbf q_l)\big),
\end{equation*}
with a two-point volume flux $\mathbf f^\#$ (e.g.\ kinetic-energy- or
entropy-conserving) and a numerical surface flux $\mathbf f^*$.
\begin{itemize}
\item \emph{Volume terms.} The two-point fluxes are products of
  arithmetic means, e.g.\ for $\rho\theta$:
  $\{\!\{\rho u\}\!\}\{\!\{\theta\}\!\}$ or
  $\{\!\{\rho\}\!\}\{\!\{u\}\!\}\{\!\{\theta\}\!\}$. Freezing all factors
  except one fast factor makes them bilinear. Since the mean is linear,
  $\sum_jp_{ij}\{\!\{\theta_j\}\!\}_{lm}=\{\!\{\sum_jp_{ij}\theta_j\}\!\}_{lm}$.
  Hence only nodal combined fields $W_{ik}$ have to be stored, not
  pairwise quantities. Option B gives $\{\!\{\pi_y(\rho\theta)_z\}\!\}$.
\item \emph{Surface terms.} As numerical flux we use the low Mach
  number approximate Riemann solver LMARS \cite{Chen2013}, which is
  common in atmospheric models. See the paragraph below.
\item \emph{Expensive nonlinear parts avoided in the inner steps:} the
  equation of state at every node, and the logarithmic means of
  entropy-conserving fluxes. Flux differencing still costs
  $O(N^{d+1})$ per element, but with a few multiplications per pair
  instead of a full nonlinear flux evaluation.
\item Entropy stability or kinetic-energy preservation of the full split
  form is in general \emph{not} inherited by the frozen operator. Whether
  a frozen split form with a discrete energy estimate for the linear
  acoustic subsystem exists is an open question.
\end{itemize}

\paragraph{The LMARS flux and its frozen form.}
We follow Chen et al.\ \cite{Chen2013}. Their model uses the same
prognostic variables as \eqref{eq:euler} (with the scaled potential
temperature density $Q=\rho\theta/p_0^{\kappa}$, $\kappa=R_d/c_p$). Let
$\mathbf q^\pm$ be the reconstructed states on both sides of an interface
(in \cite{Chen2013} by a conservative five-point interpolation), with
normal velocity $u^\pm$ and pressure $p^\pm$. Let $\bar\rho\bar c$ be
the acoustic impedance, evaluated with a locally constant sound speed
$\bar c$ (low Mach number assumption). LMARS defines
\begin{equation*}
 u^*=\tfrac12(u^-+u^+)-\frac{\beta}{2\bar\rho\bar c}(p^+-p^-),\qquad
 p^*=\tfrac12(p^-+p^+)-\frac{\beta\,\bar\rho\bar c}{2}(u^+-u^-),
\end{equation*}
with $\beta=1$ in the standard form and $\beta\le1$ in the tuned, less
diffusive form of \cite{Chen2013}. The flux is
$\mathbf f^*=u^*\,\mathbf q^{\mathrm{up}}+p^*\mathbf n$ with the upwind
state $\mathbf q^{\mathrm{up}}=\mathbf q^-$ if $u^*>0$ and $\mathbf q^+$
otherwise. For the vertical direction \cite{Chen2013} also give a
Lagrangian version with impedance $C=\sqrt{\gamma_a\rho p}$. The dissipation acts only on pressure and
velocity jumps and is scaled by the acoustic impedance. Advected
quantities ($\theta$, tracers) are only upwinded with $u^*$ and not
diffused with the sound speed. A Rusanov flux, in contrast, damps all
jumps with $\lambda\approx|u|+c$, which is far too diffusive for advected
quantities at low Mach number.

The structure of LMARS matches the splitting: the acoustic part
($u^*$, $p^*$) is fast, while the upwind selection is a slow, frozen
decision. A frozen form, consistent at $z=y$, is
\begin{align*}
 u^*_z&=\tfrac12\Big(\frac{(\rho u)^-_z}{\rho^-_y}+\frac{(\rho u)^+_z}{\rho^+_y}\Big)
       -\frac{\beta\,(p^+_z-p^-_z)}{2(\bar\rho\bar c)_y},\\
 p^*_z&=\tfrac12(p^-_z+p^+_z)-\frac{\beta\,(\bar\rho\bar c)_y}{2}
       \Big(\frac{(\rho u)^+_z}{\rho^+_y}-\frac{(\rho u)^-_z}{\rho^-_y}\Big),\\
 \text{mass flux}&=\rho^{\mathrm{up}}_y\,u^*_z,\qquad
 \rho\theta\text{ flux}=(\rho\theta)^{\mathrm{up}}_y\,u^*_z,\\
 \text{momentum flux}&=\textcolor{red}{u^*(\rho\mathbf v)^{\mathrm{up}}}(y)+p^*_z\mathbf n,
\end{align*}
with
\begin{itemize}
\item the upwind side determined by the sign of $u^*$ at the frozen
  state, and $\bar\rho\bar c$, $1/\rho^\pm$ frozen;
\item $p_z=\pi_y(\rho\theta)_z$ in Option B, or
  $p_z=p_y+a_y\big((\rho\theta)_z-(\rho\theta)_y\big)$ in Option A$'$.
\end{itemize}
Every term is a frozen coefficient times a fast variable, so the flux is
affine in $z$. Combinations $\sum_jp_{ij}(\tau)(\cdot)_{Y_j}$ of the
frozen face coefficients can be formed once per stage, as in
Lemma~\ref{lem:summation}. With $z=y$ the original LMARS flux is
recovered exactly. The same frozen LMARS flux can also be used in finite
volume schemes, where it replaces the centred pressure gradient and mass
flux of the C-grid formulation.

% =====================================================================
\section{The inner problems: structure, solution, and cost}
\label{sec:inner}

\subsection{Structure}

With Lemma~\ref{lem:summation} the inner problem of stage $i$ for a
frozen system of the form \eqref{eq:bilinear} is the linear
non-autonomous ODE
\begin{equation}\label{eq:innerlinear}
 Z_i'(\tau)=h\Big(r_i(\tau)+\mathcal B_i(\tau)Z_i\Big),\qquad
 \mathcal B_i(\tau)=q_i(\tau)\mathfrak B_0+\mathfrak B_1\big(W_i(\tau)\big),
 \qquad\tau\in[0,1],
\end{equation}
where $r_i$, $W_i$ (and $q_i$) are polynomials of degree $K$. Special
cases:
\begin{itemize}
\item $K=0$ (constant coupling): \eqref{eq:innerlinear} is autonomous
  with constant forcing.
\item \eqref{eq:L}: $\mathcal B_i(\tau)=q_i(\tau)L$. All operators
  commute, and the solution is given by $\varphi$-functions
  (Section~\ref{sec:exponential}).
\item $K\ge1$ with genuinely state-dependent $\mathfrak B_1$: the operators
  $\mathcal B_i(\tau)$ at different times do not commute.
\end{itemize}

\subsection{Overview of solution methods}

\paragraph{(a) Explicit small steps.} This is the standard in NWP:
$m_i$ steps of size $1/m_i$ in $\tau$ (i.e.\ $c_ih/m_i$ physically) with
\begin{itemize}
\item a forward--backward scheme \cite{Mesinger1977}: momentum forward,
  then $\rho,\rho\theta$ with the new momentum;
\item or a low-stage RK scheme;
\item usually combined with vertically implicit (off-centred
  Crank--Nicolson) treatment of the vertical acoustic and buoyancy terms
  (HEVI), and with divergence damping \cite{KlempSkamarockDudhia2007,Baldauf2010}.
\end{itemize}
The cost is one application of $\mathcal B_i$ per small step (plus
column solves). For the order analysis the inner method enters as in
Section~\ref{sec:nonexact}. For a linear non-autonomous problem with
polynomial coefficients, the inner RK method needs only the order
conditions for linear time-dependent problems. These lie between the
linear autonomous conditions ($b^TA^{k-1}\mathbb 1=1/k!$) and the full
nonlinear set.

\paragraph{(b) Exponential / Krylov methods (autonomous case).}
For $K=0$, or with $\mathcal B_i$ frozen over a substep, the exact
solution is
\begin{equation*}
 Z_i(1)=e^{h\mathcal B_i}Z_i(0)+\sum_{k=0}^Kk!\,h\,\varphi_{k+1}(h\mathcal B_i)\,R_{ik}.
\end{equation*}
\begin{itemize}
\item All terms can be obtained from one action of the exponential of an
  augmented matrix of dimension $N+K+1$
  \cite{Sidje1998,AlMohyHigham2011}.
\item Possible algorithms: Krylov (Arnoldi) methods \cite{NiesenWright2012,GaudreaultRainwaterTokman2018},
  polynomial interpolation at Leja points \cite{CaliariKandolfOstermannRainer2014},
  or Chebyshev expansions.
\item For (nearly) skew-adjoint acoustic operators the spectrum lies on
  the imaginary axis. A Chebyshev expansion then needs about
  $\rho(h\mathcal B_i)$ matrix--vector products, i.e.\ about the number
  of explicit small steps at CFL number one. It integrates exactly up to
  a tolerance, without CFL safety factor and without numerical damping.
\item Chebyshev and Leja methods need only spectral bounds, no inner
  products, and hence no global communication. Krylov methods need
  global reductions in every iteration, a drawback on parallel machines.
\end{itemize}

\paragraph{(c) Magnus-type methods (non-autonomous case).}
For $K\ge1$ the solution operator of \eqref{eq:innerlinear} can be
approximated by
\begin{itemize}
\item Magnus integrators \cite{BlanesCasasOteoRos2009}, e.g.\ the fourth
  order scheme with Gauss points $\tau_{1,2}$,
  \begin{equation*}
  Z_i(1)\approx e^{\Omega}Z_i(0),\qquad
  \Omega=\tfrac h2\big(\mathcal B_i(\tau_1)+\mathcal B_i(\tau_2)\big)+\tfrac{\sqrt3h^2}{12}\big[\mathcal B_i(\tau_2),\mathcal B_i(\tau_1)\big]
  \end{equation*}
  (for the homogeneous part; the forcing is added by augmentation);
\item or commutator-free Magnus integrators
  \cite{BlanesMoan2006,AlvermannFehske2011}, products of exponentials of
  linear combinations $\sum_l\beta_l\mathcal B_i(\tau_l)$.
\end{itemize}
Because $\mathfrak B_1$ is linear in $w$, a linear combination
$\sum_l\beta_l\mathfrak B_1(W_i(\tau_l))=\mathfrak B_1(\sum_l\beta_lW_i(\tau_l))$
is again a single bilinear operator, so each exponential costs the same
as in the autonomous case. A commutator-free Magnus scheme for the inner
problem is therefore itself a sequence of MIS substages with constant
coupling. Analysing this refinement within Definition~\ref{def:mis} is
an interesting consistency check.

\paragraph{(d) Implicit inner steps.}
The inner problems are linear, so one or a few implicit steps
(Crank--Nicolson, DIRK) need only linear solves with
$I-\delta h\mathcal B_i$. With a single implicit step per stage, MIS
methods become semi-implicit / IMEX schemes of the form
\cite{BoscarinoFilbetRusso2016} (Section~\ref{sec:specialcases}). With
implicit treatment only in the vertical they are HEVI schemes. The
linear operator changes in every stage, so the preconditioner (or
factorization) has to be updated per stage, not per inner step.

\begin{center}
\begin{tabular}{p{3.2cm}p{3.3cm}p{3.6cm}p{3.6cm}}
\toprule
inner solver & accuracy & cost per stage & remarks\\\midrule
explicit small steps & order $q$, CFL-limited & $m_i$ applications of $\mathcal B_i$ & local communication only; damping needed\\
Chebyshev / Leja & tolerance & $\approx\rho(h\mathcal B_i)$ applications & no inner products; autonomous case\\
Krylov & tolerance & few $\times$ Krylov dim. & global reductions\\
(CF-)Magnus & order 4--6 & several exponentials & non-autonomous, $K\ge1$\\
implicit (HEVI) & order 1--2 & linear solves & no CFL limit in vertical\\
\bottomrule
\end{tabular}
\end{center}

\subsection{How much cheaper is the frozen linearized system?}
\label{sec:cost}

Let $C_F$ be the cost of one evaluation of the full semi-discrete
right-hand side $\mathcal F_h$, and $C_B$ the cost of one inner-step
application of \eqref{eq:innerlinear}, i.e.\ of $\mathcal B_i(\tau)Z$
plus forming $W_i(\tau)$ and $r_i(\tau)$. Per large step there are $s$
stage evaluations of $f$ and of the frozen fields $w$ (cost $\approx sC_F$)
and in total $M=\sum_im_i$ inner steps. Then:
\begin{align*}
 \text{single-rate explicit with the fast step:} &\quad \approx M\,C_F,\\
 \text{MIS, naive } \textstyle\sum_jp_{ij}F(Y_j,Z): &\quad \approx sC_F+\sum_i(i-1)m_i\,C_F,\\
 \text{MIS, bilinear (Lemma~\ref{lem:summation}):} &\quad \approx sC_F+M\,C_B .
\end{align*}
For $M\gg s$ the gain over a single-rate explicit scheme is
$\approx C_F/C_B$, and the gain over the naive evaluation grows with
the number of coupled stages.

\paragraph{What is saved in $C_B$.}
Per cell and inner step the frozen system avoids:
\begin{itemize}
\item the equation of state (a \texttt{pow} per cell; frozen $a$ or $\pi$
  instead);
\item the high-order reconstructions of $u,w,\theta$ (WENO5: candidate
  stencils, smoothness indicators and nonlinear weights, of the order of
  $50$--$100$ flops per face and variable), replaced by a multiplication
  with the frozen face value;
\item momentum advection (frozen completely, part of $r_i$);
\item physics and diffusion (part of $r_i$);
\item in DG: the nonlinear two-point fluxes (logarithmic means), the
  nodal equation of state, and the impedance, upwind decision and
  velocities in the LMARS flux.
\end{itemize}
What remains is mass flux divergence, pressure gradient, buoyancy,
$K$ axpy for $W_i(\tau)$, the forcing $r_i(\tau)$ and, with HEVI,
tridiagonal column solves.

\paragraph{Rough estimate (2D, FV, WENO5).}
These are counts under stated assumptions, to be replaced by
measurements.
\begin{itemize}
\item Full right-hand side: about 6 reconstructions per cell (3
  variables, 2 directions) at 50--100 flops, plus fluxes, divergence and
  equation of state, in total $C_F\approx400$--$700$ flops per cell.
\item Frozen system: about $25$ flops per cell plus $\approx6(K+1)$ for
  forming $W_i(\tau)$ and $r_i(\tau)$, so $C_B\approx30$--$50$ flops per
  cell.
\item The flop ratio is $C_F/C_B\approx10$--$20$.
\end{itemize}
The picture is different for \emph{memory traffic}. The frozen step
reads the state and, additionally, $K+1$ copies of the frozen
coefficient fields ($\hat\theta$ on all faces, $a$ or $\pi$) and of the
forcing. This is comparable to or larger than the traffic of the full
right-hand side. On bandwidth-limited hardware (GPUs) the realistic gain
is therefore smaller, perhaps a factor $2$--$5$.
\todo{measure; consider recomputing $W_i(\tau)$ on the fly vs.\ storing.}

\paragraph{Communication.}
The inner operator has stencil width one (second-order centred pressure
gradient and divergence). WENO5 needs width three. Halo exchanges in the
inner loop are therefore narrower (less volume), although equally
frequent (same latency).

\paragraph{Options A, A$'$, B.}
The inner-step cost is the same for all three options. Option A$'$
needs per stage in addition the vector $\mathfrak A(Y_j)Y_j$, one extra
bilinear evaluation per stage, which is negligible against $M\,C_B$.

\section{Non-exact solution of the inner problems}
\label{sec:nonexact}

Section~\ref{sec:inner} lists the practical inner solvers. Here we ask
what accuracy they must have.

In practice the inner problems \eqref{eq:mis2} are solved by $m$ steps of
an explicit RK method $(A^f,b^f,c^f)$ of order $q$ with step
$\delta=1/m$ in $\tau$. Then the whole scheme is again a (generalized
additive / partitioned) RK method, and the same trees apply.

\paragraph{Computing the conditions.}
In the algorithm of Section~\ref{sec:algorithm} the exact integration
$\int_0^\tau\phi$ is replaced by $m$ steps of the inner RK method
applied to the triangular system for $\{\zeta_i(\cdot;u)\}_{u\le t}$.
This yields coefficients $\eta^{(m)}_{s+1}(t)$ for every $m$. Because
the inner equations are polynomial in $\tau$ and in the lower
coefficients, one expects an asymptotic expansion
\begin{equation*}
 \eta^{(m)}_{s+1}(t)=\eta_{s+1}(t)+\sum_{k\ge q}e_k(t)\,m^{-k},
\end{equation*}
with $e_k(t)=0$ for trees without white vertices (no inner
integration is needed for the black-only part when $\gamma$-terms and the
polynomial forcing are integrated exactly).

\paragraph{How good does the inner solution have to be?}
\begin{itemize}
\item \emph{Asymptotic view, $m$ fixed, $h\to0$.} An error term
  $e_k(t)m^{-k}h^{|t|}$ in the local error of a tree of order $|t|\le p$
  reduces the global order to $|t|-1$. Hence for order $p$ with every
  $m$ the inner method must reproduce $\eta_{s+1}(t)$ \emph{exactly}
  for all trees with $|t|\le p$. This is a finite set of algebraic
  conditions on $(A^f,b^f,c^f)$ that depends on the outer method and on
  the coupling degree $K$.
\item \emph{Conjecture.} For constant coupling ($K=0$), $g_{ij}\equiv0$, an
  inner RK method of order $q\ge p$ preserves order $p$ for every
  $m\ge1$. For $K>0$ the inner forcing is a polynomial of degree $K$ in
  $\tau$ and the inner method additionally has to integrate
  $\tau^K$-weighted terms correctly, so a higher inner order (or
  quadrature order) may be required. \todo{test computationally with the
  modified algorithm.}
\item \emph{NWP regime.} In practice the fast step $h/m$ is fixed by the
  acoustic CFL condition, so $(h/m)\|\partial_zF\|=O(1)$ and the inner
  integration is \emph{not} in its asymptotic regime. The fast acoustic
  modes are unresolved anyway. Accuracy is needed for the slow modes,
  which interact with the inner solver only through trees with white
  vertices. For the fast modes, stability (and possibly damping, e.g.\
  divergence damping, forward--backward schemes) is what matters.
\end{itemize}

\begin{question}
Find conditions on the inner method that guarantee order $p$ of the
slow component (trees with black root and at least one black vertex)
while only requiring stability for the fast modes. Can this be expressed
with the linear stability function
$R(z_s,z_f;m)$ of the test equation $y'=\lambda_sy+\lambda_fy$,
$F(y,z)=\lambda_sy+\lambda_fz$, for $z_f$ on the imaginary axis?
\end{question}

% =====================================================================
\section{Open questions and next steps}
\label{sec:open}

\begin{enumerate}
\item \textbf{Relations.} Precise embedding of MRI-GARK, MERK,
  commutator-free Lie group methods and exponential RK methods with mixed
  arguments into Definition~\ref{def:mis}; consequences for the order
  conditions (do the known conditions of these classes coincide with
  ours restricted to the respective trees?).
\item \textbf{Closed forms.} Closed-form conditions up to order 4 for
  the polynomial class; a simplifying assumption of stage order type
  \eqref{eq:Cq}; counting independent conditions.
\item \textbf{Stiff order.} Necessity in Proposition~\ref{prop:stiff};
  hyperbolic (skew-adjoint) fast operators; a nonlinear stiff theory
  for \eqref{eq:G}.
\item \textbf{Inner solution.} Order conditions with finite $m$; the
  conjecture of Section~\ref{sec:nonexact}; error expansion in $m^{-1}$.
\item \textbf{Pressure options.} Linear stability of Options A$'$ and B
  for a split acoustic--gravity test problem: does the residual acoustic
  fraction $(\gamma_a-1)/\gamma_a$ in the slow part of Option B restrict
  the large step? Is the term $-\delta\nabla_ha$ of Option A$'$ harmless?
\item \textbf{Polynomial $g_{ij}$.} Benefit of polynomial $g_{ij}$
  (Section~\ref{sec:polygamma}) for stability; additional stiff conditions.
\item \textbf{Discretizations.} A frozen split-form DG operator with an
  energy estimate for the acoustic subsystem; measured values of
  $C_F/C_B$ for FV--WENO and DG on CPUs and GPUs.
\item \textbf{Construction.} Methods of order 4 and 5 with polynomial
  coupling, with additional criteria (error constants, linear
  stability for $z_f\in i\R$, uniform fast speed, $c_i$ monotone so
  that the inner integration runs forward in time).
\end{enumerate}

% =====================================================================
% =====================================================================
\appendix
\section{Order conditions of order 4}
\label{app:order4}

Conditions $\sum_i\hat b_id_i(t)=1/\gamma(t)$ for the 26 trees of order
4, general method of Definition~\ref{def:mis}. Notation as in
Section~\ref{sec:orderconditions}; $\eta_j(u)$ and $\zeta_i(\sigma;u)$
for $|u|\le3$ are given by Theorem~\ref{thm:general} and
Table~\ref{tab:oc3}. Column ``class'' refers to the subclass
\eqref{eq:restart}: B = base method, F = automatically satisfied, C1--C4
= the coupling conditions of Proposition~\ref{prop:count}, -- = implied
by lower order, base and C1--C4.

{\renewcommand{\arraystretch}{1.5}
\begin{longtable}{rcllc}
\caption{Order 4 conditions.}\label{tab:oc4}\\
\toprule
No. & $t$ & $d_i(t)$ & $1/\gamma(t)$ & class\\\midrule
\endfirsthead
\toprule
No. & $t$ & $d_i(t)$ & $1/\gamma(t)$ & class\\\midrule
\endhead
\bottomrule
\endfoot
1 & \btree{[,b [,b [,b [,b]]]]} & $\sum_j P_{ij}(1)\,\eta_j(\Tbbbl)$ & $\frac1{24}$ & B\\[3pt]
2 & \btree{[,b [,b [,b [,w]]]]} & $\sum_j P_{ij}(1)\,\eta_j(\Tbbwl)$ & $\frac1{24}$ & --\\[3pt]
3 & \btree{[,b [,b [,b][,b]]]} & $\sum_j P_{ij}(1)\,\eta_j(\Tbbb)$ & $\frac1{12}$ & B\\[3pt]
4 & \btree{[,b [,b [,b][,w]]]} & $\sum_j P_{ij}(1)\,\eta_j(\Tbbw)$ & $\frac1{12}$ & C4\\[3pt]
5 & \btree{[,b [,b [,b]][,w]]} & $\sum_j \eta_j(\Tbb)\int_0^1 p_{ij}(\sigma)\,z_i(\sigma)\dd\sigma$ & $\frac1{8}$ & C3\\[3pt]
6 & \btree{[,b [,b [,w [,b]]]]} & $\sum_j P_{ij}(1)\,\eta_j(\Tbwbl)$ & $\frac1{24}$ & --\\[3pt]
7 & \btree{[,b [,b [,w [,w]]]]} & $\sum_j P_{ij}(1)\,\eta_j(\Tbwwl)$ & $\frac1{24}$ & --\\[3pt]
8 & \btree{[,b [,b [,w][,w]]]} & $\sum_j P_{ij}(1)\,\eta_j(\Tbww)$ & $\frac1{12}$ & --\\[3pt]
9 & \btree{[,b [,b [,w]][,w]]} & $\sum_j \eta_j(\Tbw)\int_0^1 p_{ij}(\sigma)\,z_i(\sigma)\dd\sigma$ & $\frac1{8}$ & --\\[3pt]
10 & \btree{[,b [,b][,b [,b]]]} & $\sum_j P_{ij}(1)\,c_j\eta_j(\Tbb)$ & $\frac1{8}$ & B\\[3pt]
11 & \btree{[,b [,b][,b [,w]]]} & $\sum_j P_{ij}(1)\,c_j\eta_j(\Tbw)$ & $\frac1{8}$ & --\\[3pt]
12 & \btree{[,b [,b][,b][,b]]} & $\sum_j P_{ij}(1)\,c_j^{3}$ & $\frac1{4}$ & B\\[3pt]
13 & \btree{[,b [,b][,b][,w]]} & $\sum_j c_j^{2}\int_0^1 p_{ij}(\sigma)\,z_i(\sigma)\dd\sigma$ & $\frac1{4}$ & C1\\[3pt]
14 & \btree{[,b [,b][,w [,b]]]} & $\sum_j c_j\int_0^1 p_{ij}(\sigma)\,\zeta_i(\sigma;\Tbb)\dd\sigma$ & $\frac1{8}$ & --\\[3pt]
15 & \btree{[,b [,b][,w [,w]]]} & $\sum_j c_j\int_0^1 p_{ij}(\sigma)\,\zeta_i(\sigma;\Tbw)\dd\sigma$ & $\frac1{8}$ & --\\[3pt]
16 & \btree{[,b [,b][,w][,w]]} & $\sum_j c_j\int_0^1 p_{ij}(\sigma)\,\big(z_i(\sigma)\big)^{2}\dd\sigma$ & $\frac1{4}$ & C2\\[3pt]
17 & \btree{[,b [,w [,b [,b]]]]} & $\int_0^1 q_i(\sigma)\,\zeta_i(\sigma;\Tbbbl)\dd\sigma$ & $\frac1{24}$ & --\\[3pt]
18 & \btree{[,b [,w [,b [,w]]]]} & $\int_0^1 q_i(\sigma)\,\zeta_i(\sigma;\Tbbwl)\dd\sigma$ & $\frac1{24}$ & --\\[3pt]
19 & \btree{[,b [,w [,b][,b]]]} & $\int_0^1 q_i(\sigma)\,\zeta_i(\sigma;\Tbbb)\dd\sigma$ & $\frac1{12}$ & --\\[3pt]
20 & \btree{[,b [,w [,b][,w]]]} & $\int_0^1 q_i(\sigma)\,\zeta_i(\sigma;\Tbbw)\dd\sigma$ & $\frac1{12}$ & --\\[3pt]
21 & \btree{[,b [,w [,w [,b]]]]} & $\int_0^1 q_i(\sigma)\,\zeta_i(\sigma;\Tbwbl)\dd\sigma$ & $\frac1{24}$ & --\\[3pt]
22 & \btree{[,b [,w [,w [,w]]]]} & $\int_0^1 q_i(\sigma)\,\zeta_i(\sigma;\Tbwwl)\dd\sigma$ & $\frac1{24}$ & F\\[3pt]
23 & \btree{[,b [,w [,w][,w]]]} & $\int_0^1 q_i(\sigma)\,\zeta_i(\sigma;\Tbww)\dd\sigma$ & $\frac1{12}$ & F\\[3pt]
24 & \btree{[,b [,w][,w [,b]]]} & $\int_0^1 q_i(\sigma)\,z_i(\sigma)\zeta_i(\sigma;\Tbb)\dd\sigma$ & $\frac1{8}$ & --\\[3pt]
25 & \btree{[,b [,w][,w [,w]]]} & $\int_0^1 q_i(\sigma)\,z_i(\sigma)\zeta_i(\sigma;\Tbw)\dd\sigma$ & $\frac1{8}$ & F\\[3pt]
26 & \btree{[,b [,w][,w][,w]]} & $\int_0^1 q_i(\sigma)\,\big(z_i(\sigma)\big)^{3}\dd\sigma$ & $\frac1{4}$ & F\\[3pt]
\end{longtable}}

\begin{thebibliography}{99}
\small
\bibitem{AlMohyHigham2011} A.\,H.~Al-Mohy, N.\,J.~Higham, Computing the
 action of the matrix exponential, with an application to exponential
 integrators, \emph{SIAM J. Sci. Comput.} 33 (2011) 488--511.
\bibitem{AlvermannFehske2011} A.~Alvermann, H.~Fehske, High-order
 commutator-free exponential time-propagation of driven quantum systems,
 \emph{J. Comput. Phys.} 230 (2011) 5930--5956.
\bibitem{Baldauf2010} M.~Baldauf, Linear stability analysis of
 Runge--Kutta-based partial time-splitting schemes for the Euler
 equations, \emph{Mon. Wea. Rev.} 138 (2010) 4475--4496.
\bibitem{BlanesCasasOteoRos2009} S.~Blanes, F.~Casas, J.\,A.~Oteo,
 J.~Ros, The Magnus expansion and some of its applications, \emph{Phys.
 Rep.} 470 (2009) 151--238.
\bibitem{BlanesMoan2006} S.~Blanes, P.\,C.~Moan, Fourth- and sixth-order
 commutator-free Magnus integrators for linear and non-linear dynamical
 systems, \emph{Appl. Numer. Math.} 56 (2006) 1519--1537.
\bibitem{CaliariKandolfOstermannRainer2014} M.~Caliari, P.~Kandolf,
 A.~Ostermann, S.~Rainer, Comparison of software for computing the action
 of the matrix exponential, \emph{BIT} 54 (2014) 113--128.
\bibitem{GassnerWintersKopriva2016} G.\,J.~Gassner, A.\,R.~Winters,
 D.\,A.~Kopriva, Split form nodal discontinuous Galerkin schemes with
 summation-by-parts property for the compressible Euler equations,
 \emph{J. Comput. Phys.} 327 (2016) 39--66.
\bibitem{GaudreaultRainwaterTokman2018} S.~Gaudreault, G.~Rainwater,
 M.~Tokman, KIOPS: a fast adaptive Krylov subspace solver for exponential
 integrators, \emph{J. Comput. Phys.} 372 (2018) 236--255.
\bibitem{GiraldoKellyConstantinescu2013} F.\,X.~Giraldo, J.\,F.~Kelly,
 E.\,M.~Constantinescu, Implicit-explicit formulations of a
 three-dimensional nonhydrostatic unified model of the atmosphere (NUMA),
 \emph{SIAM J. Sci. Comput.} 35 (2013) B1162--B1194.
\bibitem{JiangShu1996} G.-S.~Jiang, C.-W.~Shu, Efficient implementation
 of weighted ENO schemes, \emph{J. Comput. Phys.} 126 (1996) 202--228.
\bibitem{KlempSkamarockDudhia2007} J.\,B.~Klemp, W.\,C.~Skamarock,
 J.~Dudhia, Conservative split-explicit time integration methods for the
 compressible nonhydrostatic equations, \emph{Mon. Wea. Rev.} 135 (2007)
 2897--2913.
\bibitem{Mesinger1977} F.~Mesinger, Forward-backward scheme, and its use
 in a limited area model, \emph{Contrib. Atmos. Phys.} 50 (1977) 200--210.
\bibitem{NiesenWright2012} J.~Niesen, W.\,M.~Wright, Algorithm 919: a
 Krylov subspace algorithm for evaluating the $\varphi$-functions appearing
 in exponential integrators, \emph{ACM Trans. Math. Softw.} 38 (2012) 22.
\bibitem{Sidje1998} R.\,B.~Sidje, Expokit: a software package for
 computing matrix exponentials, \emph{ACM Trans. Math. Softw.} 24 (1998)
 130--156.
\bibitem{BauerKnoth2021} S.~Bauer, O.~Knoth,
 Extended multirate infinitesimal step methods: derivation of order
 conditions, \emph{J. Comput. Appl. Math.} 387 (2021) 112541.
 \todo{verify}
\bibitem{BoscarinoFilbetRusso2016} S.~Boscarino, F.~Filbet, G.~Russo,
 High order semi-implicit schemes for time dependent partial differential
 equations, \emph{J. Sci. Comput.} 68 (2016) 975--1001.
\bibitem{Butcher2016} J.\,C.~Butcher, \emph{Numerical Methods for
 Ordinary Differential Equations}, 3rd ed., Wiley, 2016.
\bibitem{CelledoniMarthinsenOwren2003} E.~Celledoni, A.~Marthinsen,
 B.~Owren, Commutator-free Lie group methods, \emph{Future Gener. Comput.
 Syst.} 19 (2003) 341--352.
\bibitem{Chen2013} X.~Chen, N.~Andronova, B.~van Leer, J.\,E.~Penner,
 J.\,P.~Boyd, C.~Jablonowski, S.-J.~Lin, A control-volume model of the
 compressible Euler equations with a vertical Lagrangian coordinate,
 \emph{Mon. Wea. Rev.} 141 (2013) 2526--2544, doi:10.1175/MWR-D-12-00129.1.
\bibitem{ChinomonaReynolds2021} R.~Chinomona, D.\,R.~Reynolds,
 Implicit-explicit multirate infinitesimal GARK methods, \emph{SIAM J.
 Sci. Comput.} 43 (2021) A3082--A3113.
\bibitem{CoxMatthews2002} S.\,M.~Cox, P.\,C.~Matthews, Exponential time
 differencing for stiff systems, \emph{J. Comput. Phys.} 176 (2002)
 430--455.
\bibitem{GuntherSandu2016} M.~G\"unther, A.~Sandu, Multirate generalized
 additive Runge--Kutta methods, \emph{Numer. Math.} 133 (2016) 497--524.
\bibitem{HairerNorsettWanner1993} E.~Hairer, S.\,P.~N{\o}rsett,
 G.~Wanner, \emph{Solving Ordinary Differential Equations I}, 2nd ed.,
 Springer, 1993.
\bibitem{HochbruckOstermann2005} M.~Hochbruck, A.~Ostermann, Explicit
 exponential Runge--Kutta methods for semilinear parabolic problems,
 \emph{SIAM J. Numer. Anal.} 43 (2005) 1069--1090.
\bibitem{HochbruckOstermann2010} M.~Hochbruck, A.~Ostermann, Exponential
 integrators, \emph{Acta Numer.} 19 (2010) 209--286.
\bibitem{HochbruckOstermannSchweitzer2009} M.~Hochbruck, A.~Ostermann,
 J.~Schweitzer, Exponential Rosenbrock-type methods, \emph{SIAM J. Numer.
 Anal.} 47 (2009) 786--803.
\bibitem{KetchesonRanocha2023} D.\,I.~Ketcheson, H.~Ranocha, Computing
 with B-series, \emph{ACM Trans. Math. Softw.} 49 (2023) 14.
\bibitem{KlempWilhelmson1978} J.\,B.~Klemp, R.\,B.~Wilhelmson, The
 simulation of three-dimensional convective storm dynamics, \emph{J.
 Atmos. Sci.} 35 (1978) 1070--1096.
\bibitem{KnothWensch2014} O.~Knoth, J.~Wensch, Generalized split-explicit
 Runge--Kutta methods for the compressible Euler equations, \emph{Mon.
 Wea. Rev.} 142 (2014) 2067--2081.
\bibitem{KnothWolke1998} O.~Knoth, R.~Wolke, Implicit-explicit
 Runge--Kutta methods for computing atmospheric reactive flows,
 \emph{Appl. Numer. Math.} 28 (1998) 327--341.
\bibitem{Lawson1967} J.\,D.~Lawson, Generalized Runge--Kutta processes
 for stable systems with large Lipschitz constants, \emph{SIAM J. Numer.
 Anal.} 4 (1967) 372--380.
\bibitem{LuanChinomonaReynolds2020} V.\,T.~Luan, R.~Chinomona,
 D.\,R.~Reynolds, A new class of high-order methods for multirate
 differential equations, \emph{SIAM J. Sci. Comput.} 42 (2020)
 A1245--A1268.
\bibitem{Owren2006} B.~Owren, Order conditions for commutator-free Lie
 group methods, \emph{J. Phys. A} 39 (2006) 5585--5599.
\bibitem{Sandu2019} A.~Sandu, A class of multirate infinitesimal GARK
 methods, \emph{SIAM J. Numer. Anal.} 57 (2019) 2300--2327.
\bibitem{SchlegelKnothArnoldWolke2009} M.~Schlegel, O.~Knoth, M.~Arnold,
 R.~Wolke, Multirate Runge--Kutta schemes for advection equations,
 \emph{J. Comput. Appl. Math.} 226 (2009) 345--357.
\bibitem{WenschKnothGalant2009} J.~Wensch, O.~Knoth, A.~Galant, Multirate
 infinitesimal step methods for atmospheric flow simulation, \emph{BIT}
 49 (2009) 449--473.
\bibitem{WickerSkamarock2002} L.\,J.~Wicker, W.\,C.~Skamarock,
 Time-splitting methods for elastic models using forward time schemes,
 \emph{Mon. Wea. Rev.} 130 (2002) 2088--2097.
\end{thebibliography}
\noindent\todo{verify all bibliographic data.}

\end{document}

